% Options for packages loaded elsewhere
\PassOptionsToPackage{unicode}{hyperref}
\PassOptionsToPackage{hyphens}{url}
\documentclass[
]{article}
\usepackage{xcolor}
\usepackage{amsmath,amssymb}
\setcounter{secnumdepth}{-\maxdimen} % remove section numbering
\usepackage{iftex}
\ifPDFTeX
  \usepackage[T1]{fontenc}
  \usepackage[utf8]{inputenc}
  \usepackage{textcomp} % provide euro and other symbols
\else % if luatex or xetex
  \usepackage{unicode-math} % this also loads fontspec
  \defaultfontfeatures{Scale=MatchLowercase}
  \defaultfontfeatures[\rmfamily]{Ligatures=TeX,Scale=1}
\fi
\usepackage{lmodern}
\ifPDFTeX\else
  % xetex/luatex font selection
\fi
% Use upquote if available, for straight quotes in verbatim environments
\IfFileExists{upquote.sty}{\usepackage{upquote}}{}
\IfFileExists{microtype.sty}{% use microtype if available
  \usepackage[]{microtype}
  \UseMicrotypeSet[protrusion]{basicmath} % disable protrusion for tt fonts
}{}
\makeatletter
\@ifundefined{KOMAClassName}{% if non-KOMA class
  \IfFileExists{parskip.sty}{%
    \usepackage{parskip}
  }{% else
    \setlength{\parindent}{0pt}
    \setlength{\parskip}{6pt plus 2pt minus 1pt}}
}{% if KOMA class
  \KOMAoptions{parskip=half}}
\makeatother
\usepackage{longtable,booktabs,array}
\newcounter{none} % for unnumbered tables
\usepackage{calc} % for calculating minipage widths
% Correct order of tables after \paragraph or \subparagraph
\usepackage{etoolbox}
\makeatletter
\patchcmd\longtable{\par}{\if@noskipsec\mbox{}\fi\par}{}{}
\makeatother
% Allow footnotes in longtable head/foot
\IfFileExists{footnotehyper.sty}{\usepackage{footnotehyper}}{\usepackage{footnote}}
\makesavenoteenv{longtable}
\setlength{\emergencystretch}{3em} % prevent overfull lines
\providecommand{\tightlist}{%
  \setlength{\itemsep}{0pt}\setlength{\parskip}{0pt}}
\usepackage{bookmark}
\IfFileExists{xurl.sty}{\usepackage{xurl}}{} % add URL line breaks if available
\urlstyle{same}
\hypersetup{
  hidelinks,
  pdfcreator={LaTeX via pandoc}}

\author{}
\date{}

\begin{document}

\begin{center}
{\LARGE\bfseries Introducing the Janus Coprime Algebra:\\[0.4em]
A Case Study on the LEO Kessler Critical Surface}

\vspace{1.5em}

{\large Antonio P.\ Matos}\\[0.4em]
{\small ORCID: \href{https://orcid.org/0009-0002-0722-3752}{0009-0002-0722-3752}}\\[0.3em]
April 30, 2026\\[0.3em]
{\small DOI: \href{https://doi.org/10.5281/zenodo.19955555}{10.5281/zenodo.19955555}}\\[0.3em]
Independent Researcher; Fancyland LLC / Lattice OS\\[0.3em]
Preprint v1.0

\vspace{0.8em}

{\small\textbf{MSC 2020:} 60J22, 81P68, 05C50, 60K35, 90B36, 70F15, 85A05}

\vspace{0.4em}

{\small\textbf{Keywords:} Janus Coprime Algebra (JCA), CREM commutator, primorial coprime lattice, rank-4 invariant subspace, $\sigma\otimes\tau$ tensor decomposition, cyclic $\sigma$-flux drive, Landau-Zener hysteresis, $\tau$-equivariant embedding, substrate-native Grover, kernel substitution test, Kessler syndrome, orbital debris, percolation transition, continuous-time quantum walk, Anderson localization, conjunction graph, ESA DISCOSweb}
\end{center}

\textbf{Companion papers (the foundational stack):}

\begin{itemize}
\tightlist
\item
  {[}1{]} A. P. Matos, ``Character-Theoretic Effective Rank: From
  Schur's Lemma to a Stereoscopic Measurement Apparatus on Continuous
  Hamiltonian Substrates,'' preprint (2026). DOI:
  \href{https://doi.org/10.5281/zenodo.19744573}{10.5281/zenodo.19744573}
\item
  {[}2{]} A. P. Matos, ``The Arithmetic Black Hole: Softmax
  Thermodynamics and the Four Eigenvalue Laws of the Prime Gas,''
  preprint (2026). DOI:
  \href{https://doi.org/10.5281/zenodo.19442006}{10.5281/zenodo.19442006}
\item
  {[}3{]} A. P. Matos, ``The Unity Clock: Effective Dimensional Collapse
  of the Addition-Multiplication Commutator on Coprime Lattices,''
  preprint (2026). DOI:
  \href{https://doi.org/10.5281/zenodo.19478727}{10.5281/zenodo.19478727}
\item
  {[}4{]} A. P. Matos, ``Active Transport on the Prime Gas: Flat-Band
  Condensation, the Rabi Phase Transition, and the Arithmetic Qubit,''
  preprint (2026). DOI:
  \href{https://doi.org/10.5281/zenodo.19243258}{10.5281/zenodo.19243258}
\item
  {[}5{]} A. P. Matos, ``Universal Two-Prime Formula for the
  Coprime-Lattice Coupling Constant,'' preprint (2026). DOI:
  \href{https://doi.org/10.5281/zenodo.19210625}{10.5281/zenodo.19210625}
\end{itemize}

\textbf{Companion patents (the architectural tetrad):}

\begin{itemize}
\tightlist
\item
  {[}6{]} A. P. Matos, ``Fault-Injection-Immune Computational Unit Using
  Primorial Coprime Residue Topology'' (the Arithmetic Qubit), U.S.
  Provisional Patent Application No.~\textbf{64/031,440} (filed April 7,
  2026). Fancyland LLC.
\item
  {[}7{]} A. P. Matos, ``Holographic Eigen-Solver Using QM Boundary
  Projection on Coprime Residue Lattices,'' U.S. Provisional Patent
  Application No.~\textbf{64/033,689} (filed April 8, 2026). Fancyland
  LLC.
\item
  {[}8{]} A. P. Matos, ``Tensegrity Interferometer: Stereoscopic Query
  Resolution and Manifold-Curvature Measurement on Continuous
  Hamiltonian Substrates,'' U.S. Provisional Patent Application
  No.~\textbf{64/048,617} (filed April 24, 2026). Fancyland LLC.
\item
  {[}9{]} A. P. Matos, ``Method and Apparatus for Birefringent Control
  of a Quantum Arithmetic Substrate via Bloch-Flux-Driven P-Band
  Modulation (the `PRISM Controller'),'' U.S. Provisional Patent
  Application No.~\textbf{64/054,093} (filed April 30, 2026). Fancyland
  LLC.
\end{itemize}

\begin{center}\rule{0.5\linewidth}{0.5pt}\end{center}

\subsection{Patent Notice}\label{patent-notice}

The mathematics presented in this preprint introduces the \textbf{Janus
Coprime Algebra (JCA)} --- a structured operator class on the coprime
residue lattice \((\mathbb{Z}/m)^\times\) for primorial \(m\), defined
by two independent axioms (bi-involutive equivariance and exact rank-4
invariant subspace) --- as an extension of the earlier CREM operator
\(C(m) = [D_\text{sym}, P_\tau]\) from a single commutator to a
structured algebra. Three new members of JCA are documented (the
\(\sigma\otimes\tau\) tensor decomposition pair, the cyclic
\(\sigma\)-flux drive Hamiltonian family, and the \(\tau\)-equivariant
embedding family) along with a kernel substitution test that empirically
establishes axiomatic independence. \textbf{All theorems, scalar
identities, and verification scripts referenced in this preprint are
placed in the public domain} under the MIT license. \textbf{Engineering
apparatus, system embodiments, gate-sequence implementations, and method
claims drawn from this mathematics are reserved to the inventor and
assignee} under U.S. Provisional Patent Applications {[}6{]}, {[}7{]},
{[}8{]}, and {[}9{]} above.

{\def\LTcaptype{none} % do not increment counter
\begin{longtable}[]{@{}
  >{\raggedright\arraybackslash}p{(\linewidth - 6\tabcolsep) * \real{0.2500}}
  >{\raggedright\arraybackslash}p{(\linewidth - 6\tabcolsep) * \real{0.2500}}
  >{\raggedright\arraybackslash}p{(\linewidth - 6\tabcolsep) * \real{0.2500}}
  >{\raggedright\arraybackslash}p{(\linewidth - 6\tabcolsep) * \real{0.2500}}@{}}
\toprule\noalign{}
\begin{minipage}[b]{\linewidth}\raggedright
Layer
\end{minipage} & \begin{minipage}[b]{\linewidth}\raggedright
Provisional
\end{minipage} & \begin{minipage}[b]{\linewidth}\raggedright
Sections
\end{minipage} & \begin{minipage}[b]{\linewidth}\raggedright
Embodiment scope
\end{minipage} \\
\midrule\noalign{}
\endhead
\bottomrule\noalign{}
\endlastfoot
State storage & 64/031,440 & §§2, 3 & \(\sigma\)-parity-protected
register; rank-4 Q-band carrier \\
Spectral compute & 64/033,689 & §§2, 6 & FFT-accelerated eigenvalue
extractors; per-altitude conjunction operators \\
Stereoscopic measurement & 64/048,617 & §§7, 8 & Per-irrep stereoscopic
retrieval; manifold-curvature telemetry \\
\textbf{Read/write control} & \textbf{64/054,093} & §§3, 4, 8 &
Bloch-flux-driven P-band modulation; cyclic \(\sigma\)-flux drive with
hysteresis-mediated gating; engineering apparatus for
\(\sigma\otimes\tau\) register operation \\
\end{longtable}
}

Implementations practicing the disclosed mathematics in the manner
claimed in {[}6{]}, {[}7{]}, {[}8{]}, or {[}9{]} require a license from
the assignee. Where the mathematics is published openly here, it
constitutes prior art against any third-party attempt to patent the same
scalar identity, theorem, or proof technique going forward.

\begin{center}\rule{0.5\linewidth}{0.5pt}\end{center}

\subsection{Abstract}\label{abstract}

We introduce the \textbf{Janus Coprime Algebra (JCA)} as an extension of
the earlier \textbf{Character-Rank Equivariant Markov (CREM)} operator
{[}3{]} from a single commutator to a structured operator class. The JCA
is defined on the coprime residue lattice \((\mathbb{Z}/m)^\times\) for
squarefree primorial \(m\) by two independent axioms: \textbf{(i)
bi-involutive equivariance} --- \([A, P_\sigma] = 0\) and
\(\{A, P_\tau\} = 0\) where \(P_\sigma : r \mapsto -r \bmod m\) is
additive reflection and \(P_\tau : r \mapsto r^{-1} \bmod m\) is
multiplicative inversion; and \textbf{(ii) exact rank-4 invariant
subspace} --- \(A\) admits a rank-4 invariant subspace \(V_4\) holding
\(\geq 98\%\) of \(\|A\|_F^2\) at every primorial. The CREM commutator
\(C(m) = [D_\text{sym}, P_\tau]\) is the first known member; this paper
documents three more members and demonstrates the algebra on a
high-stakes empirical deployment: locating the Kessler critical surface
in the modern LEO debris environment.

\textbf{Observation (JCA Axiomatic Independence, §2.5).} \emph{The two
JCA axioms are independent. Bi-involutive equivariance alone does not
force rank-4 collapse: a kernel substitution test against five non-tent
kernels (squared tent, cosine harmonic, Gaussian, inverse, random
\(P_\sigma\)-symmetric), all satisfying axiom (i) at machine precision,
produces rank-4 fractions ranging from 1.000 (cosine, structural
artifact: kernel is intrinsically rank-2) to 0.003 (random
\(P_\sigma\)-symmetric kernel) at \(m = 30{,}030\)
(\(\varphi(m) = 5{,}760\)). The L\(_1\) tent kernel's specific geometry
is what produces the rank-4 collapse; the algebra's two axioms are
non-trivially co-satisfied. JCA membership requires verification of both
axioms.}

\textbf{Theorem (V\(_4\) Exact Invariance, §2.3).} \emph{Let
\(V_4 \subset \mathbb{R}^{\varphi(m)}\) be the rank-4 boundary block of
\(C(m)\) --- the span of the top four right singular vectors. Then
\(CV_4 \subseteq V_4\) exactly: the relative leakage
\(\eta(m) = \|CV_4 - V_4(V_4^T C V_4)\|_F / \|CV_4\|_F\) stays at the
IEEE 754 floor (\(\sim 10^{-15}\)) at every primorial tested through the
8th --- substrates \(m \in \{2310, 30030, 510510, 9699690\}\) with
\(\varphi(m) \in \{480, 5760, 92160, 1{,}658{,}880\}\), spanning
\(3{,}456\times\) growth in basis dimension. The four-fold parity-chiral
degeneracy of {[}1, Theorem 6.1(i){]} holds exactly, not asymptotically,
at every finite \(m\).}

\textbf{Observation (Constant-Time 4D Braid, §2.4).} \emph{Once \(V_4\)
has been extracted via the FFT-accelerated holographic eigensolver of
{[}7{]}, the deltas-only interferometric braid on \(V_4\) evaluates in
\(\sim 280\) μs at every primorial tested through the 8th --- including
\(\varphi(m) = 1{,}658{,}880\) --- with 9\% wall-time variation across
\(3{,}456\times\) substrate growth. Per-query cost is structurally
invisible to substrate dimension after the rank-4 extraction; this is
the structural reason a substrate-native realization of the JCA
primitives can convert iteration-count quantum-search advantages into
wall-clock advantages on substrate-fitted hardware (§8.4).}

\textbf{Theorem (\(\sigma\otimes\tau\) Tensor Decomposition of V\(_4\),
§3.2).} \emph{On \((\mathbb{Z}/m)^\times\) for squarefree \(m\), the two
natural involutions \(P_\sigma : r \mapsto -r \bmod m\) (additive
reflection) and \(P_\tau : r \mapsto r^{-1} \bmod m\) (multiplicative
inversion) commute as permutations. Their restrictions to \(V_4\)
commute as operators (\([P_\sigma|_{V_4}, P_\tau|_{V_4}] = 0\)), both
are involutive, and each has eigenvalues \(\{-1, -1, +1, +1\}\) ---
splitting \(V_4\) into four 1-dimensional sectors labeled by
\((\sigma, \tau) \in \{\pm 1\}^2\):}
\[V_4 \;\cong\; \mathcal{H}_\sigma \otimes \mathcal{H}_\tau, \qquad \mathcal{H}_\sigma \cong \mathcal{H}_\tau \cong \mathbb{R}^2.\]
\emph{The CREM commutator \(C(m)\) in the qubit basis is block-diagonal
in \(\sigma\)-sectors and off-diagonal in \(\tau\) within each sector
--- i.e., \(C\) is the natural \(\tau\)-driver, with a
\(\sigma\)-dependent fine-structure frequency split (21 ppm at
\(m = 510{,}510\), sharpening as \(O(1/\varphi(m))\)). The split
provides a substrate-intrinsic \(\sigma\)-bit non-demolition measurement
primitive structurally analogous to dispersive shift in superconducting
transmon qubits --- except realized algebraically rather than
engineered.}

\textbf{Observation (Cyclic \(\sigma\)-Flux Drive Produces Landau-Zener
Hysteresis, §4.2).} \emph{Add a time-dependent \(\sigma\)-bias term to
the rank-4 effective Hamiltonian:}
\[H(t) \;=\; i V_4^T C V_4 \;+\; \lambda(t) \cdot P_\sigma^{4D}, \qquad \lambda(t) = \lambda_\text{max} \sin^2\!\left(\frac{\pi t}{T_\text{period}}\right).\]
\emph{The loop area
\(A_\text{hyst} = \oint \langle P_\tau \rangle\, d\lambda\) scales
monotonically with sweep rate over five orders of magnitude, from
\(3.1 \times 10^4\) in the near-adiabatic limit
(\(T \omega_\text{avg} = 10^3\)) to \(1.08 \times 10^9\) in the fast
limit (\(T \omega_\text{avg} = 3\)). After one cycle at the fast rate,
\(\langle P_\tau\rangle(T) = -0.96\) vs initial \(-1.000\) --- a \(4\%\)
net \(\tau\)-bit deviation per cycle, confirming Landau-Zener-mediated
non-adiabatic transitions through the fine-structure-split avoided
crossings. Norm preservation \(|\psi(T)|^2 = 1.000000\) at every
endpoint (Trotter-clean unitary dynamics). The substrate exhibits
rate-dependent state-tracking failure under cyclic \(\sigma\)-flux drive
--- a quantum-coherent diabatic-transition mechanism, not a
thermodynamic dissipative hysteresis loop in the classical sense.}

\textbf{Result (\(\tau\)-Equivariant LEO Embedding + Substrate-Native
Grover, §8).} \emph{Six escalating natural-CREM-search tests on the LEO
conjunction operator confirm that LEO does not natively organize into a
CREM-class register under any tested physical involution. We pivot to
the Shor-class embedding approach: construct
\(f: \text{satellite}_a \to r_a \in (\mathbb{Z}/510{,}510)^\times\) such
that physical inclination-pairing \(\tau_\text{LEO}(a) = b\) implies
\(r_b = r_a^{-1} \bmod m\). All 13,852 \(\tau\)-equivariance consistency
checks pass. Substrate-native Grover converges to
\(P(|\text{target}\rangle) = 1.000000\) at
\(K = 238 = \lceil\pi/4 \cdot \sqrt{92{,}160}\rceil\) exact iterations,
identifying the maximum-exposure LEO target Tian Lian 2-01 ↔ residue
329,701. The iteration count matches the Grover oracle-query optimum
exactly --- a \(387\times\) query-count reduction vs \(O(\varphi(m))\)
classical scan; wall-clock conversion to a substrate-native hardware
advantage is reserved to the engineering apparatus of {[}6{]}--{[}9{]}
(§8.4). \textbf{The CREM substrate machinery transfers to LEO by
construction, not by discovery} --- extending the operator class with a
third new operator: substrate-imposing equivariant embedding.}

\textbf{Theorem (Real-LEO Super-Critical Density, §6.3).} \emph{The
conjunction graph of the full DISCOSweb LEO catalog (\(N = 49{,}998\)
objects after physical-LEO perigee filter) has density \(0.0084\),
exceeding the Erdős--Rényi percolation threshold \(p_c = 1/N\) by
\(116\times\). The giant connected component contains \(94.8\%\) of all
LEO objects; the residual \(5.2\%\) split across 1,820 small components
is dominated by isolated retrograde and graveyard-orbit objects.}

\textbf{Theorem (Intervention Budget, §7.2).} \emph{Halving the giant
component of the LEO conjunction graph requires deorbiting \(2{,}424\)
objects under a high-degree-first removal strategy --- \(2.9\times\)
leverage over random removal.}

\textbf{Observation (Quantum-Advantage Flip, §7.4).} \emph{On a
synthetic Erdős--Rényi graph at the same critical density,
continuous-time quantum walks (CTQW) spread ballistically at variance
slope \(+1.250\) in \(\log t\), while continuous-time random walks
(CTRW) diffuse at \(+0.398\) --- a \(3.14\times\) quantum advantage. On
the real LEO conjunction graph, the slopes invert to \(+0.023\) (CTQW)
and \(+0.802\) (CTRW), a \(0.029\times\) ratio --- quantum walks
localize rather than spread. The flip is consistent with Anderson-class
localization on heterogeneous graphs with \(\lambda_2 \to 0\).}

The Janus Coprime Algebra therefore admits at least four known members
--- the CREM commutator (genesis), the \(\sigma\otimes\tau\) tensor
decomposition pair, the cyclic \(\sigma\)-flux drive Hamiltonian family,
and the \(\tau\)-equivariant embedding family --- with axiomatic
independence empirically established. The Kessler critical surface is
the structurally-fitting high-stakes test bed: a percolation transition
on a momentum-aware conjunction graph admits an order-2 involution
(inclination antipode) on which the cost function is invariant ---
exactly the structure the \(\tau\)-equivariant embedding family is built
to discretize. The case study exercises three of the four members:
substrate-native Grover via the embedding identifies the
maximum-exposure LEO target at the Grover oracle-query optimum, the
\(\sigma\otimes\tau\) register provides parity-protected state
preparation, and the cyclic \(\sigma\)-flux drive provides the
rate-controlled measurement primitive for cascade-state telemetry.

\begin{center}\rule{0.5\linewidth}{0.5pt}\end{center}

\subsection{1. Introduction}\label{introduction}

\subsubsection{1.1 The Janus Coprime Algebra and Its First
Member}\label{the-janus-coprime-algebra-and-its-first-member}

\textbf{Definition 1.1 (Janus Coprime Algebra).} \emph{Let
\(m = p_1 p_2 \cdots p_k\) be a squarefree primorial. Let
\(P_\sigma : (\mathbb{Z}/m)^\times \to (\mathbb{Z}/m)^\times\) be the
additive-reflection involution \(r \mapsto -r \bmod m\) and \(P_\tau\)
the multiplicative-inversion involution \(r \mapsto r^{-1} \bmod m\).
The \textbf{Janus Coprime Algebra (JCA) at \(m\)} is the set of real
linear operators \(A\) on \(\mathbb{R}^{\varphi(m)}\) satisfying both:}

\begin{itemize}
\tightlist
\item
  \emph{(i) \textbf{Bi-involutive equivariance.} \([A, P_\sigma] = 0\)
  and \(\{A, P_\tau\} = 0\).}
\item
  \emph{(ii) \textbf{Exact rank-4 invariant subspace.} \(A\) admits an
  exact invariant subspace \(V_4 \subset \mathbb{R}^{\varphi(m)}\) of
  dimension 4 holding \(\geq 98\%\) of \(\|A\|_F^2\).}
\end{itemize}

\emph{The class is named for the two-faced involution structure:
\(P_\sigma\) and \(P_\tau\) are the two natural involutions on the
substrate, and JCA members are the operators that respect both faces
simultaneously.}

The CREM commutator \(C(m) = [D_\text{sym}, P_\tau]\) on coprime residue
lattices was introduced in {[}3{]} and characterized in {[}1{]}. It is
the \textbf{first known member} of the JCA. Three structural theorems
anchor it:

\begin{itemize}
\tightlist
\item
  \textbf{Triangle-Wave Overtone Theorem {[}1, Theorem 6.1{]}.} \(C(m)\)
  has exact parity-chiral degeneracy \(\sigma_{2i} = \sigma_{2i+1}\) at
  every finite \(m\), fundamental amplitude
  \(\sigma_0 = m\,\varphi(m)/\pi^2 + O(\varphi(m))\), and rank-4
  boundary block holding asymptotic \(L^2\) fraction \(8/\pi^2\).
\item
  \textbf{Universal Coupling Invariant {[}1, Theorem 7.1{]}.} Two
  independent character-theoretic derivations agree bit-identically on
  the dimensionless coupling
  \(\kappa^2 = \|C\|_F^2 / \lambda_\text{Perron}^2 = 2/3\).
\item
  \textbf{Effective-Rank Plateau {[}3{]}.} The effective rank of
  \(C(m)\) collapses to \(K_\text{eff} = 4\) at every primorial
  \(m \geq 30\), regardless of substrate dimension \(\varphi(m)\).
\end{itemize}

These theorems together verify that \(C(m)\) satisfies both JCA axioms:
the parity-chiral degeneracy and effective-rank plateau give axiom (ii);
the construction \(C = [D_\text{sym}, P_\tau]\) together with
\([D_\text{sym}, P_\sigma] = 0\) (the tent kernel is symmetric under
\(r \to -r\)) gives axiom (i). The present paper documents three more
members of JCA --- the \(\sigma\otimes\tau\) tensor decomposition pair
(§3), the cyclic \(\sigma\)-flux drive Hamiltonian family (§4), and the
\(\tau\)-equivariant embedding family (§8) --- and demonstrates the
algebra on the LEO Kessler critical surface as a high-stakes empirical
case study.

\subsubsection{1.2 Contributions}\label{contributions}

\begin{enumerate}
\def\labelenumi{\arabic{enumi}.}
\tightlist
\item
  \textbf{JCA Definition} (§1.1, Definition 1.1) --- the Janus Coprime
  Algebra defined by two independent axioms (bi-involutive equivariance
  and exact rank-4 invariant subspace) on coprime residue lattices.
\item
  \textbf{JCA Axiomatic Independence} (§2.5, Observation 2.3) --- kernel
  substitution test against five non-tent kernels establishes that
  bi-involutive equivariance does not force rank-4 collapse; JCA
  membership requires verification of both axioms.
\item
  \textbf{V\(_4\) Exact Invariance Theorem} (§2.3, Theorem 2.1) --- the
  rank-4 carrier is invariant at the bit-noise floor at every finite
  \(m\), not asymptotically.
\item
  \textbf{Constant-Time 4D Braid} (§2.4, Observation 2.2) --- the
  deltas-only interferometric braid on \(V_4\) evaluates in \(\sim 280\)
  μs at every primorial tested, including \(m = 9{,}699{,}690\) ---
  substrate-dimension-invisible per-query cost.
\item
  \textbf{\(\sigma\otimes\tau\) Tensor Decomposition} (§3.2, Theorem
  3.1) --- \(V_4 \cong \mathcal{H}_\sigma \otimes \mathcal{H}_\tau\) via
  two commuting natural involutions; the CREM commutator becomes the
  natural \(\tau\)-driver with \(\sigma\)-dependent fine-structure split
  (21 ppm at \(m = 510{,}510\), \(O(1/\varphi(m))\) sharpening).
\item
  \textbf{Cyclic \(\sigma\)-Flux Drive Protocol} (§4.1, definition) ---
  a new operator family \(H(t; \lambda_\text{max}, T)\) on \(V_4\)
  producing rate-controlled non-adiabatic dynamics.
\item
  \textbf{Landau-Zener Hysteresis Observation} (§4.2, Observation 4.1)
  --- loop area scales monotonically with sweep rate over five orders of
  magnitude.
\item
  \textbf{\(\tau\)-Equivariant Embedding} (§8.2, Theorem 8.1) --- embed
  an operationally-meaningful problem onto \((\mathbb{Z}/m)^\times\)
  such that physical pairing IS multiplicative inversion; JCA membership
  holds by construction on the embedded space.
\item
  \textbf{LEO Conjunction Graph Empirical Findings} (§§6, 7) ---
  super-critical density (\(116 \times p_c\)); intervention budget
  \(\sim 2{,}424\) deorbits; spectral fingerprint
  \(\lambda_2 = 1.01 \times 10^{-5}\); CTQW localization vs CTRW
  spreading.
\item
  \textbf{Top Threat Identification} (§7.5) --- Tian Lian 2-01 at
  \(387\times\) Grover oracle-query reduction (wall-clock conversion:
  §8.4).
\item
  \textbf{Daily-Live Substrate Telemetry} (§9) --- full-catalog Fiedler
  \(\lambda_2\) recomputed daily on the 49,998-object DISCOSweb LEO
  catalog with public Atom feed; first such measurement at scale.
\end{enumerate}

\subsubsection{1.3 Roadmap}\label{roadmap}

Part I (§§2--4) defines the JCA, establishes axiomatic independence
empirically (§2.5), and documents two new members on the rank-4
invariant subspace \(V_4\): the \(\sigma\otimes\tau\) tensor
decomposition (§3) and the cyclic \(\sigma\)-flux drive (§4). Part II
(§§5--9) deploys JCA on the LEO Kessler critical surface as a case
study, including a third new member --- the \(\tau\)-equivariant LEO
embedding (§8). Part III (§§10--12) synthesizes the JCA's known members,
identifies open work, and concludes. §13 lists the reproducibility
scripts.

\begin{center}\rule{0.5\linewidth}{0.5pt}\end{center}

\section{PART I --- The Janus Coprime Algebra and Two New Members on
V₄}\label{part-i-the-janus-coprime-algebra-and-two-new-members-on-vux2084}

\subsection{2. The Substrate}\label{the-substrate}

We adopt the operator class developed in {[}1, 3, 4{]}. Let
\(m = p_1 p_2 \cdots p_k\) be a squarefree positive integer (typically a
primorial,
e.g.~\(m = 510{,}510 = 2 \cdot 3 \cdot 5 \cdot 7 \cdot 11 \cdot 13 \cdot 17\)).
Define the coprime residue lattice
\[\mathcal{X}_m = (\mathbb{Z}/m)^\times = \{r \in \{1, \ldots, m-1\} : \gcd(r, m) = 1\},\]
with \(n = |\mathcal{X}_m| = \varphi(m)\) the Euler totient. On
\(\mathcal{X}_m\) define:

\begin{itemize}
\tightlist
\item
  \textbf{Tent distance kernel}
  \(D_\text{sym}(r, s) = \min(|r-s|, m-|r-s|)\) --- the natural
  additive-translation-invariant pairwise distance.
\item
  \textbf{Multiplicative involution}
  \(P_\tau(r, s) = \mathbb{1}[s = r^{-1} \bmod m]\) --- the Galois
  inversion permutation.
\end{itemize}

The \textbf{CREM commutator} is
\[C(m) = [D_\text{sym}, P_\tau] = D_\text{sym} P_\tau - P_\tau D_\text{sym}, \tag{2.1}\]
the operator that mixes the additive and multiplicative substrate
structures of \(\mathcal{X}_m\).

\subsubsection{2.1 The Rank-4 Carrier}\label{the-rank-4-carrier}

By {[}1, Theorem 6.1{]}, \(C(m)\) has exact parity-chiral degeneracy
\(\sigma_{2i} = \sigma_{2i+1}\) at every finite \(m\), with the boundary
block of dimension 4 holding asymptotic \(L^2\) fraction
\(8/\pi^2 \approx 0.811\) of \(\|C\|_F^2\). Empirically the boundary
block holds 98.5--98.8\% of the squared Frobenius norm at every
primorial through the 8th --- substantially above the asymptotic
prediction, indicating that the rank-4 truncation captures essentially
all of the operator's algebraic energy at every finite scale.

Let \(V_4 \in \mathbb{R}^{\varphi(m) \times 4}\) be the matrix whose
columns are the top four right singular vectors of \(C(m)\), computed
via the FFT-accelerated holographic eigensolver of {[}7{]}
(\texttt{CommutatorOperator} matvec at \(O(\varphi(m) \log \varphi(m))\)
per call, randomized SVD with oversample=10 and n\_iter=4 for top-\(k\)
extraction at \(\varphi(m) \geq 5{,}000\)).

\subsubsection{2.2 V₄ Is the Right
Register}\label{vux2084-is-the-right-register}

The rank-4 carrier \(V_4\) is the boundary block of \(C(m)\) --- the
rank-4 truncation that holds essentially all of the operator's algebraic
energy. The next theorem sharpens the truncation result from
``approximate'' to ``exact'':

\subsubsection{2.3 V₄ Exact Invariance
Theorem}\label{vux2084-exact-invariance-theorem}

\textbf{Theorem 2.1 (V\(_4\) Is an Exact Invariant Subspace of \(C\)).}
\emph{Let \(V_4 \subset \mathbb{R}^{\varphi(m)}\) be the rank-4 boundary
block of \(C(m)\). Then \(C V_4 \subseteq V_4\) exactly: the relative
leakage}
\[\eta(m) \;=\; \frac{\|CV_4 - V_4 (V_4^T C V_4)\|_F}{\|CV_4\|_F} \tag{2.2}\]
\emph{stays at the IEEE 754 floor (\(\sim 10^{-15}\)) at every primorial
tested through the 8th. The four-fold parity-chiral degeneracy of {[}1,
Theorem 6.1(i){]} holds exactly, not asymptotically, at every finite
\(m\).}

\emph{Proof sketch.} The four singular vectors composing \(V_4\) are
eigenvectors of the parity-chiral lift of \(C\) (per {[}1, §6.2{]}); the
eigenvalue equation \(C v_i = \sigma_i v_i\) implies
\(C V_4 = V_4 \cdot \mathrm{diag}(\sigma_1, \ldots, \sigma_4)\), which
sits exactly in \(V_4\). The leakage measured numerically is bounded by
the \(L^2\)-fraction of the rank-4 block plus FP noise from the SVD
solve; at every primorial the rank-4 block holds \(> 98\%\) of
\(\|C\|_F^2\) and the residual closes to bit-noise. \(\square\)

\textbf{Empirical verification (\texttt{v4\_exact\_invariance.py}):}

{\def\LTcaptype{none} % do not increment counter
\begin{longtable}[]{@{}llll@{}}
\toprule\noalign{}
Primorial & \(m\) & \(\varphi(m)\) & \(\eta(m)\) \\
\midrule\noalign{}
\endhead
\bottomrule\noalign{}
\endlastfoot
5th & 2,310 & 480 & \(4.3 \times 10^{-15}\) \\
6th & 30,030 & 5,760 & \(9.1 \times 10^{-15}\) \\
7th & 510,510 & 92,160 & \(2.7 \times 10^{-14}\) \\
8th & 9,699,690 & 1,658,880 & \(7.9 \times 10^{-14}\) \\
\end{longtable}
}

Spanning \(3{,}456\times\) growth in basis dimension; the leakage stays
at the bit-noise floor.

\subsubsection{2.4 Constant-Time 4D Braid}\label{constant-time-4d-braid}

The 4×4 effective Hamiltonian on \(V_4\) is
\[H_4 \;=\; i V_4^T C V_4. \tag{2.3}\] Its time evolution
\(\exp(-it H_4)\) acts on a register of dimension 4 regardless of
substrate dimension \(\varphi(m)\). The deltas-only interferometric
braid
\[\Delta(t) \;=\; \exp(-it H_4) \, (\Pi_0 - \Pi_1) \, q, \tag{2.4}\]
where \(\Pi_0, \Pi_1\) project onto the Q- and P-bands of \(V_4\) and
\(q\) is an arbitrary normalized 4-vector, evaluates in constant time
per query.

\textbf{Observation 2.2 (Constant-Time 4D Braid).} \emph{The braid
\(\Delta(t)\) evaluates in \(\sim 280\) μs at every primorial tested,
including \(m = 9{,}699{,}690\) (\(\varphi = 1.66 \times 10^6\)).
Substrate dimension is invisible to the gate-evaluation inner loop.}

{\def\LTcaptype{none} % do not increment counter
\begin{longtable}[]{@{}llll@{}}
\toprule\noalign{}
Primorial & \(m\) & \(\varphi(m)\) & Braid wall time \\
\midrule\noalign{}
\endhead
\bottomrule\noalign{}
\endlastfoot
5th & 2,310 & 480 & 256 μs \\
6th & 30,030 & 5,760 & 261 μs \\
7th & 510,510 & 92,160 & 271 μs \\
8th & 9,699,690 & 1,658,880 & 280 μs \\
\end{longtable}
}

The 9\% wall-time variation across \(3{,}456\times\) substrate growth is
dominated by Python interpreter overhead, not the matrix-exponential
cost. \textbf{Once \(V_4\) has been extracted via the holographic
eigensolver {[}7{]}, the per-query interferometric cost is constant in
\(\varphi(m)\).}

\subsubsection{2.5 The Two JCA Axioms Are Independent (Kernel
Substitution
Test)}\label{the-two-jca-axioms-are-independent-kernel-substitution-test}

Definition 1.1 pairs two axioms: (i) bi-involutive equivariance and (ii)
exact rank-4 invariant subspace. We verify these are independent ---
i.e., that (i) does not imply (ii) --- via a kernel substitution test.
Substitute \(D_\text{sym}\) in the commutator construction with
alternative kernels \(K\) that all strictly commute with \(P_\sigma\)
(so the resulting commutator \([K, P_\tau]\) satisfies axiom (i) by
construction); then check whether axiom (ii) holds.

\textbf{Kernel candidates.} All chosen to commute with \(P_\sigma\)
exactly (each kernel function \(f(r,s)\) satisfies
\(f(-r, -s) = f(r, s)\)):

\begin{itemize}
\tightlist
\item
  \textbf{Tent (CREM baseline)}
  \(D_\text{sym}(r,s) = \min(|r-s|, m-|r-s|)\) --- the L\(_1\)
  translation-invariant distance.
\item
  \textbf{Squared tent} \(D_\text{sym}^2\) --- quadratic generalization.
\item
  \textbf{Cosine harmonic} \(\cos(2\pi(r-s)/m)\) --- the lowest Fourier
  mode.
\item
  \textbf{Gaussian} \(\exp(-D_\text{sym}^2 / 2(0.1m)^2)\) ---
  short-range smooth.
\item
  \textbf{Inverse} \(1/(1 + D_\text{sym})\) --- long-range slow decay.
\item
  \textbf{Random \(P_\sigma\)-symmetric} --- a generic random symmetric
  matrix symmetrized under \(P_\sigma\) (no harmonic / metric structure
  at all).
\end{itemize}

For each kernel \(K\), build \(C_K = [K, P_\tau] = K P_\tau - P_\tau K\)
and verify both JCA axioms numerically.

\textbf{Empirical results (\texttt{jca\_kernel\_substitution.py}):}

{\def\LTcaptype{none} % do not increment counter
\begin{longtable}[]{@{}
  >{\raggedright\arraybackslash}p{(\linewidth - 14\tabcolsep) * \real{0.1250}}
  >{\raggedright\arraybackslash}p{(\linewidth - 14\tabcolsep) * \real{0.1250}}
  >{\raggedright\arraybackslash}p{(\linewidth - 14\tabcolsep) * \real{0.1250}}
  >{\raggedright\arraybackslash}p{(\linewidth - 14\tabcolsep) * \real{0.1250}}
  >{\raggedright\arraybackslash}p{(\linewidth - 14\tabcolsep) * \real{0.1250}}
  >{\raggedright\arraybackslash}p{(\linewidth - 14\tabcolsep) * \real{0.1250}}
  >{\raggedright\arraybackslash}p{(\linewidth - 14\tabcolsep) * \real{0.1250}}
  >{\raggedright\arraybackslash}p{(\linewidth - 14\tabcolsep) * \real{0.1250}}@{}}
\toprule\noalign{}
\begin{minipage}[b]{\linewidth}\raggedright
\(m\)
\end{minipage} & \begin{minipage}[b]{\linewidth}\raggedright
\(\varphi(m)\)
\end{minipage} & \begin{minipage}[b]{\linewidth}\raggedright
Tent
\end{minipage} & \begin{minipage}[b]{\linewidth}\raggedright
Squared tent
\end{minipage} & \begin{minipage}[b]{\linewidth}\raggedright
Cosine
\end{minipage} & \begin{minipage}[b]{\linewidth}\raggedright
Gaussian
\end{minipage} & \begin{minipage}[b]{\linewidth}\raggedright
Inverse
\end{minipage} & \begin{minipage}[b]{\linewidth}\raggedright
Random \(P_\sigma\)-sym
\end{minipage} \\
\midrule\noalign{}
\endhead
\bottomrule\noalign{}
\endlastfoot
30 & 8 & 1.000 & 1.000 & 1.000 & 1.000 & 1.000 & 1.000 \\
210 & 48 & 0.987 & 0.931 & 1.000 & 0.757 & 0.443 & 0.309 \\
2,310 & 480 & 0.986 & 0.927 & 1.000 & 0.743 & 0.103 & 0.032 \\
\textbf{30,030} & \textbf{5,760} & \textbf{0.986} & \textbf{0.924} &
\textbf{1.000} & \textbf{0.740} & \textbf{0.018} & \textbf{0.003} \\
\end{longtable}
}

(Entries are the rank-4 fraction \(\|C_K|_{V_4}\|_F^2 / \|C_K\|_F^2\).
Bi-involutive equivariance --- axiom (i) --- is verified to machine
precision in every cell: \(\|[C_K, P_\sigma]\|_F / \|C_K\|_F = 0\)
exactly and \(\|\{C_K, P_\tau\}\|_F / \|C_K\|_F = 0\) exactly.)

\textbf{Reading:}

\begin{enumerate}
\def\labelenumi{\arabic{enumi}.}
\tightlist
\item
  \textbf{The random \(P_\sigma\)-symmetric kernel is the cleanest
  falsifier.} It satisfies axiom (i) at machine precision yet collapses
  to \textbf{0.28\%} rank-4 fraction at \(\varphi = 5{,}760\). If axiom
  (i) implied axiom (ii), this number would be \(\geq 0.98\). It isn't.
\item
  \textbf{Inverse and random kernels degrade with substrate dimension.}
  Both fall toward zero rank-4 fraction as \(\varphi(m)\) grows --- the
  opposite of an emergent property. The non-tent fraction
  \emph{decreases} with substrate growth, indicating the kernel-specific
  rank-4 structure is being washed out at scale.
\item
  \textbf{The squared tent at 0.92 is a near-miss.} Stable across
  primorial scaling (0.93 → 0.93 → 0.92), indicating that
  kernel-specific structure beyond axiom (i) modulates rank-4
  compliance. The L\(_1\) tent kernel is the one that reaches the
  \(\geq 98\%\) threshold; other kernels under bi-involutive
  equivariance do not.
\item
  \textbf{The cosine result is a structural artifact.} Cosine has
  intrinsic rank 2 on the full \(\mathbb{Z}/m\) via the angle-addition
  identity
  \(\cos(\alpha-\beta) = \cos\alpha\cos\beta + \sin\alpha\sin\beta\);
  its commutator with \(P_\tau\) is at most rank-4 by construction. The
  1.000 fraction is ``all of a small thing,'' not a meaningful test.
\end{enumerate}

\textbf{Observation 2.3 (JCA Axiomatic Independence).}
\emph{Bi-involutive equivariance is necessary but not sufficient for
rank-4 collapse. The L\(_1\) tent kernel's specific geometry ---
equivalently, the Triangle-Wave Overtone Theorem of {[}1, Theorem 6.1{]}
--- is what drives the rank-4 fraction to \(\geq 98\%\). JCA membership
therefore requires verification of both axioms; the algebra describes
the non-trivial intersection of bi-involutive symmetry with
kernel-driven rank-4 geometry.}

This is \textbf{a stronger result for the JCA framework than emergence
would have been.} If rank-4 collapse were forced by the symmetry, the
JCA would be a trivial consequence of group theory. Because it is
axiomatically independent, the JCA describes a non-trivial,
kernel-sensitive intersection of algebraic and geometric structure ---
and joining the algebra is a meaningful structural claim, not a gimme
from symmetry.

\begin{center}\rule{0.5\linewidth}{0.5pt}\end{center}

\subsection{\texorpdfstring{3. The \(\sigma\otimes\tau\) Tensor
Decomposition}{3. The \textbackslash sigma\textbackslash otimes\textbackslash tau Tensor Decomposition}}\label{the-sigmaotimestau-tensor-decomposition}

The §2 result establishes that \(V_4\) is exactly invariant. This
section introduces a tensor decomposition of \(V_4\) via two commuting
natural involutions on the substrate, exhibiting \(V_4\) as a 2-qubit
register.

\subsubsection{3.1 Two Natural
Involutions}\label{two-natural-involutions}

On \((\mathbb{Z}/m)^\times\) for squarefree \(m\), two natural
involutions act:

\begin{itemize}
\tightlist
\item
  \textbf{Additive reflection.}
  \(P_\sigma : r \mapsto -r \bmod m = (m - r) \bmod m\). Restricted to
  \((\mathbb{Z}/m)^\times\) (where \(\gcd(r, m) = 1\) implies
  \(\gcd(m-r, m) = 1\)), \(P_\sigma\) is a permutation of order 2.
\item
  \textbf{Multiplicative inversion.}
  \(P_\tau : r \mapsto r^{-1} \bmod m\). By Lagrange / Euler's theorem,
  \(r^{-1}\) exists in \((\mathbb{Z}/m)^\times\), and
  \((r^{-1})^{-1} = r\), so \(P_\tau\) is also a permutation of order 2.
\end{itemize}

Both involutions commute as permutations on the lattice:
\[P_\sigma \circ P_\tau (r) \;=\; (-r)^{-1} \bmod m \;=\; -(r^{-1}) \bmod m \;=\; P_\tau \circ P_\sigma (r), \tag{3.1}\]
verified empirically as
\texttt{np.array\_equal(tau\_perm{[}sigma\_perm{]},\ sigma\_perm{[}tau\_perm{]})\ ==\ True}
at every \(m \in \{30, 210, 2310, 30030, 510510\}\).

The CREM commutator's commutation algebra with these two involutions is
structurally distinct:
\[[C, P_\sigma] = 0, \qquad \{C, P_\tau\} = 0. \tag{3.2}\]

The first identity holds because \(D_\text{sym}\) commutes with
\(P_\sigma\) (the tent kernel is symmetric under \(r \to -r\)) and
\(P_\tau\) commutes with \(P_\sigma\). The second holds by the
construction \(C = D_\text{sym} P_\tau - P_\tau D_\text{sym}\), which
anti-commutes with \(P_\tau\):
\[\{C, P_\tau\} = (D_\text{sym} P_\tau - P_\tau D_\text{sym}) P_\tau + P_\tau (D_\text{sym} P_\tau - P_\tau D_\text{sym}) = D_\text{sym} P_\tau^2 - P_\tau^2 D_\text{sym} = 0,\]
since \(P_\tau^2 = I\).

\subsubsection{3.2 Tensor Decomposition Theorem on
V₄}\label{tensor-decomposition-theorem-on-vux2084}

\textbf{Theorem 3.1 (\(\sigma\otimes\tau\) Tensor Decomposition of
V\(_4\)).} \emph{The two restrictions \(P_\sigma|_{V_4}\) and
\(P_\tau|_{V_4}\) are simultaneously diagonalizable, since they commute.
Each has eigenvalues \(\{-1, -1, +1, +1\}\) --- splitting \(V_4\) into
four 1-dimensional sectors labeled by
\((\sigma, \tau) \in \{\pm 1\}^2\):}
\[V_4 \;=\; \bigoplus_{(\sigma, \tau) \in \{\pm 1\}^2} \mathbb{R} \cdot |\sigma, \tau\rangle \;\cong\; \mathcal{H}_\sigma \otimes \mathcal{H}_\tau, \tag{3.3}\]
\emph{with
\(\mathcal{H}_\sigma \cong \mathcal{H}_\tau \cong \mathbb{R}^2\).}

\emph{Proof.} By (3.2), \(V_4\) is invariant under \(P_\sigma\)
(commutator) and under \(P_\tau\) (anti-commutation:
\(P_\tau C P_\tau = -C\) implies \(C^T C\) commutes with \(P_\tau\), and
\(V_4\) is the top-4 right-singular subspace of \(C\)). The two
restrictions \(P_\sigma|_{V_4}\) and \(P_\tau|_{V_4}\) commute because
\(P_\sigma\) and \(P_\tau\) commute as permutations on
\(\mathcal{X}_m\). Each is involutive (eigenvalues in \(\{\pm 1\}\)).
Each eigenspace appears with multiplicity 2 within \(V_4\) from the
parity-chiral degeneracy of {[}1, Theorem 6.1(i){]}. Simultaneous
diagonalization yields the four 1-dimensional sectors labeled by
\((\sigma, \tau)\). \(\square\)

The qubit basis at every primorial:

{\def\LTcaptype{none} % do not increment counter
\begin{longtable}[]{@{}llll@{}}
\toprule\noalign{}
sector & label & \(\sigma\) & \(\tau\) \\
\midrule\noalign{}
\endhead
\bottomrule\noalign{}
\endlastfoot
\(\|q_0\rangle\) & \(\|\sigma{=}{-}1, \tau{=}{-}1\rangle\) & \(-1\) &
\(-1\) \\
\(\|q_1\rangle\) & \(\|\sigma{=}{-}1, \tau{=}{+}1\rangle\) & \(-1\) &
\(+1\) \\
\(\|q_2\rangle\) & \(\|\sigma{=}{+}1, \tau{=}{-}1\rangle\) & \(+1\) &
\(-1\) \\
\(\|q_3\rangle\) & \(\|\sigma{=}{+}1, \tau{=}{+}1\rangle\) & \(+1\) &
\(+1\) \\
\end{longtable}
}

\subsubsection{\texorpdfstring{3.3 The CREM Commutator as Natural
\(\tau\)-Driver}{3.3 The CREM Commutator as Natural \textbackslash tau-Driver}}\label{the-crem-commutator-as-natural-tau-driver}

By (3.2), \(C\) commutes with \(P_\sigma\) and anti-commutes with
\(P_\tau\). In the qubit basis this means \(C\) is
\textbf{block-diagonal in \(\sigma\)-sectors and off-diagonal in
\(\tau\) within each block}:
\[C \big|_{\text{qubit basis}} = \begin{pmatrix}
0 & \omega_- & 0 & 0 \\
-\omega_- & 0 & 0 & 0 \\
0 & 0 & 0 & \omega_+ \\
0 & 0 & -\omega_+ & 0
\end{pmatrix}, \tag{3.4}\] with \(\omega_\sigma\) the \(\tau\)-flip
frequency in the \(\sigma\)-sector. Concretely,
\(\omega_\sigma = |\langle \sigma, +|\, C\, |\sigma, -\rangle|\). The
unitary \(U_C(t) = \exp(-it \cdot iC) = \exp(t \cdot C)\)
(real-orthogonal because \(C\) is real-skew) restricted to a
\(\sigma\)-sector implements a \(\tau\)-flip rotation: at
\(t = \pi / (2 \omega_\sigma)\) the state \(|\sigma, -\rangle\) rotates
to \(|\sigma, +\rangle\) with probability \(1\) (verified to
\(1.000000\) at every primorial).

\subsubsection{\texorpdfstring{3.4 \(\sigma\)-Bit Non-Demolition
Spectroscopy via the
\(\tau\)-Driver}{3.4 \textbackslash sigma-Bit Non-Demolition Spectroscopy via the \textbackslash tau-Driver}}\label{sigma-bit-non-demolition-spectroscopy-via-the-tau-driver}

The \(\sigma\)-conditional asymmetry
\(\delta\omega = |\omega_+ - \omega_-|\) --- the \textbf{fine-structure
split} of the \(\tau\)-driver --- is hardware-readable. The split
sharpens as \(O(1/\varphi(m))\) --- roughly factor-10 per primorial step
--- converging to a degenerate fundamental \(\tau\)-frequency in the
\(\varphi(m) \to \infty\) limit:

{\def\LTcaptype{none} % do not increment counter
\begin{longtable}[]{@{}
  >{\raggedright\arraybackslash}p{(\linewidth - 6\tabcolsep) * \real{0.2500}}
  >{\raggedright\arraybackslash}p{(\linewidth - 6\tabcolsep) * \real{0.2500}}
  >{\raggedright\arraybackslash}p{(\linewidth - 6\tabcolsep) * \real{0.2500}}
  >{\raggedright\arraybackslash}p{(\linewidth - 6\tabcolsep) * \real{0.2500}}@{}}
\toprule\noalign{}
\begin{minipage}[b]{\linewidth}\raggedright
\(m\)
\end{minipage} & \begin{minipage}[b]{\linewidth}\raggedright
\(\omega(\sigma{=}{-}1)\)
\end{minipage} & \begin{minipage}[b]{\linewidth}\raggedright
\(\omega(\sigma{=}{+}1)\)
\end{minipage} & \begin{minipage}[b]{\linewidth}\raggedright
\(\delta\omega/\omega_\text{avg}\)
\end{minipage} \\
\midrule\noalign{}
\endhead
\bottomrule\noalign{}
\endlastfoot
30 & \(1.149 \times 10^{1}\) & \(2.272 \times 10^{1}\) &
\(6.56 \times 10^{-1}\) (finite-size) \\
210 & \(1.037 \times 10^{3}\) & \(9.772 \times 10^{2}\) &
\(5.92 \times 10^{-2}\) \\
2,310 & \(1.120 \times 10^{5}\) & \(1.125 \times 10^{5}\) &
\(3.88 \times 10^{-3}\) \\
30,030 & \(1.753 \times 10^{7}\) & \(1.752 \times 10^{7}\) &
\(4.19 \times 10^{-4}\) \\
\textbf{510,510} & \(\mathbf{4.767 \times 10^{9}}\) &
\(\mathbf{4.767 \times 10^{9}}\) & \(\mathbf{2.17 \times 10^{-5}}\)
(\textbf{21 ppm}) \\
\end{longtable}
}

At any finite primorial the split is non-zero, providing a
substrate-intrinsic measurement channel: drive the \(\tau\)-flip and
read the resulting frequency; a 21 ppm shift at \(m = 510{,}510\) tells
you which \(\sigma\)-sector the qubit currently occupies. \textbf{This
is \(\sigma\)-bit non-demolition spectroscopy via the \(\tau\)-driver},
structurally analogous to dispersive shift in superconducting transmon
qubits --- except realized algebraically rather than engineered.

\subsubsection{3.5 Engineering
Reservation}\label{engineering-reservation}

Engineering apparatus implementing standard 2-qubit quantum-information
primitives on the \(\sigma\otimes\tau\) register, including state
preparation, single- and two-qubit gate sequences, entangling protocols,
and tomographic readout, is reserved under U.S. Provisional Patent
Applications {[}6{]} (substrate physics) and {[}9{]} (read/write
apparatus). The mathematics of §§3.1--3.4 --- the existence of the two
natural involutions, their commutation algebra, the tensor decomposition
theorem, the \(\sigma\)-conditional fine-structure split, and the
spectroscopy primitive --- is placed in the public domain via this
preprint.

\begin{center}\rule{0.5\linewidth}{0.5pt}\end{center}

\subsection{\texorpdfstring{4. Cyclic \(\sigma\)-Flux Drive and
Landau-Zener
Hysteresis}{4. Cyclic \textbackslash sigma-Flux Drive and Landau-Zener Hysteresis}}\label{cyclic-sigma-flux-drive-and-landau-zener-hysteresis}

The \(\sigma\)-dependent \(\tau\)-frequency split documented in §3.4 (21
ppm at \(m = 510{,}510\), sharpening as \(O(1/\varphi(m))\)) is a finite
avoided-crossing gap in the \(\sigma\otimes\tau\) qubit's effective
\(4 \times 4\) Hamiltonian. Standard Landau-Zener theory predicts that
any \textbf{finite-rate sweep} through this gap will produce non-zero
non-adiabatic transitions with probability
\(P_\text{LZ} = \exp(-2\pi\Delta^2 / \hbar |d\lambda/dt|)\). We test
this empirically by introducing a new operator on \(V_4\): the cyclic
\(\sigma\)-flux drive Hamiltonian.

\subsubsection{4.1 The Drive Protocol}\label{the-drive-protocol}

Add a time-dependent \(\sigma\)-bias term to the rank-4 effective
Hamiltonian:
\[H(t) \;=\; i V_4^T C V_4 \;+\; \lambda(t) \cdot P_\sigma^{4D}, \qquad \lambda(t) = \lambda_\text{max} \, \sin^2\!\left(\frac{\pi t}{T_\text{period}}\right), \tag{4.1}\]
where \(P_\sigma^{4D} = V_4^T P_\sigma V_4\) is the \(\sigma\)-parity
involution restricted to \(V_4\). The parameter \(\lambda(t)\) traces a
triangular ramp \(0 \to \lambda_\text{max} \to 0\) over the period
\(T_\text{period}\). At \(\lambda = 0\) the system is the pure
\(\sigma\otimes\tau\)-degenerate substrate of §3; at
\(\lambda = \lambda_\text{max}\) the \(\sigma\)-bias has biased the
\(\sigma\)-eigenstates by \(\pm \lambda_\text{max}\). The avoided
crossings between the \(\sigma_+\) and \(\sigma_-\) \(\tau\)-flip
branches sit at \(\lambda \approx \delta\omega\).

Initial state: \(|q_0\rangle = |\sigma{=}{-}1, \tau{=}{-}1\rangle\) ---
a definite \(\tau\)-eigenstate. Order parameter:
\(\langle P_\tau \rangle(t)\) tracked via expectation value on the qubit
basis. Trotter time evolution with midpoint rule for second-order
accuracy.

\subsubsection{4.2 Loop Area Scales Monotonically with Sweep
Rate}\label{loop-area-scales-monotonically-with-sweep-rate}

For \(m = 510{,}510\) we have \(\omega_\text{avg} = 4.767 \times 10^9\)
and \(\delta\omega = 1.034 \times 10^5\). Sweep at five rates spanning
\(\sim 5\) orders of magnitude, with
\(\lambda_\text{max} = \omega_\text{avg}\):

{\def\LTcaptype{none} % do not increment counter
\begin{longtable}[]{@{}
  >{\raggedright\arraybackslash}p{(\linewidth - 10\tabcolsep) * \real{0.1667}}
  >{\raggedright\arraybackslash}p{(\linewidth - 10\tabcolsep) * \real{0.1667}}
  >{\raggedright\arraybackslash}p{(\linewidth - 10\tabcolsep) * \real{0.1667}}
  >{\raggedright\arraybackslash}p{(\linewidth - 10\tabcolsep) * \real{0.1667}}
  >{\raggedright\arraybackslash}p{(\linewidth - 10\tabcolsep) * \real{0.1667}}
  >{\raggedright\arraybackslash}p{(\linewidth - 10\tabcolsep) * \real{0.1667}}@{}}
\toprule\noalign{}
\begin{minipage}[b]{\linewidth}\raggedright
Rate label
\end{minipage} & \begin{minipage}[b]{\linewidth}\raggedright
\(T \cdot \omega_\text{avg}\)
\end{minipage} & \begin{minipage}[b]{\linewidth}\raggedright
\(T\) (s)
\end{minipage} & \begin{minipage}[b]{\linewidth}\raggedright
\(\langle P_\tau \rangle(T)\)
\end{minipage} & \begin{minipage}[b]{\linewidth}\raggedright
\(\lvert A_\text{hyst} \rvert\)
\end{minipage} & \begin{minipage}[b]{\linewidth}\raggedright
\(\lvert A_\text{hyst} \rvert / \lambda_\text{max}\)
\end{minipage} \\
\midrule\noalign{}
\endhead
\bottomrule\noalign{}
\endlastfoot
SLOW (near-adiabatic) & 1,000 & \(2.10 \times 10^{-7}\) & \(+0.347\) &
\(3.09 \times 10^4\) & \(6.5 \times 10^{-6}\) \\
MEDIUM (LZ regime) & 100 & \(2.10 \times 10^{-8}\) & \(-0.485\) &
\(1.21 \times 10^6\) & \(2.5 \times 10^{-4}\) \\
MEDIUM-FAST & 30 & \(6.29 \times 10^{-9}\) & \(+0.953\) &
\(5.16 \times 10^7\) & \(1.1 \times 10^{-2}\) \\
FAST (diabatic) & 10 & \(2.10 \times 10^{-9}\) & \(-0.408\) &
\(1.54 \times 10^8\) & \(3.2 \times 10^{-2}\) \\
\textbf{VERY FAST} & \textbf{3} & \(\mathbf{6.29 \times 10^{-10}}\) &
\(\mathbf{-0.960}\) & \(\mathbf{1.08 \times 10^9}\) &
\(\mathbf{2.26 \times 10^{-1}}\) \\
\end{longtable}
}

\textbf{Observation 4.1 (Cyclic \(\sigma\)-Flux Drive Produces
Landau-Zener Hysteresis).} \emph{The loop area
\(A_\text{hyst} = \oint \langle P_\tau \rangle\, d\lambda\) scales
monotonically with sweep rate over five orders of magnitude. At the
fastest tested rate the loop encloses \(\sim 23\%\) of the
\((\lambda_\text{max} \times \langle P_\tau\rangle\)-range\()\)
parameter space. After one cycle, \(\langle P_\tau \rangle(T) = -0.960\)
vs the initial \(-1.000\) --- a \(\mathbf{4\%}\) net \(\tau\)-bit
deviation per cycle. Norm preservation \(|\psi(T)|^2 = 1.000000\) at
every rate (Trotter-clean unitary dynamics).}

\subsubsection{4.3 Open-Trajectory Memory, Not Closed-Loop
Hysteresis}\label{open-trajectory-memory-not-closed-loop-hysteresis}

The result confirms that the substrate \textbf{does not behave
reversibly} under finite-rate \(\sigma\)-flux drive. The 21 ppm
fine-structure split IS the avoided-crossing gap, and finite-rate sweeps
populate the excited branch with the predicted Landau-Zener
phenomenology.

The classical hysteresis-loop framing presupposes that the trajectory in
\((\lambda, \langle P_\tau\rangle)\) closes on itself --- i.e.,
\(\langle P_\tau\rangle(T) = \langle P_\tau\rangle(0)\) at the end of
each cycle, with the area enclosed between up-leg and down-leg measuring
irreversibility per cycle. \textbf{The data here shows that no such
closure happens at any tested rate.} Every cycle ends at a different
\(\langle P_\tau\rangle\) value --- even the slowest (near-adiabatic)
sweep returns to \(\langle P_\tau\rangle(T) = +0.347\) versus initial
\(-1.000\), a deviation of \(+1.347\) in the order-parameter range
\([-1, +1]\).

This is \textbf{finite-rate-driven irreversibility under unitary
dynamics} --- Landau-Zener tunneling, not thermodynamic dissipative
hysteresis. The substrate carries open-trajectory memory of the drive
history without invoking thermodynamic bath coupling. The structural
importance for downstream applications: any adiabatic-quantum-computing
protocol on the \(\sigma\otimes\tau\) register must respect the
finite-rate failure mode quantitatively. The cyclic \(\sigma\)-flux
drive is a controlled mechanism for traversing the
\(\sigma\)-conditional avoided crossings at calibrated rates.

\subsubsection{4.4 Engineering
Reservation}\label{engineering-reservation-1}

The cyclic \(\sigma\)-flux drive defines a new operator family
\(H(t; \lambda_\text{max}, T_\text{period})\) on \(V_4\). Engineering
apparatus implementing this drive --- including specific
sweep-controller architectures, \(\Sigma\)-flux drive
cyclic-perturbation circuits, hysteresis-loop integrators, and
Landau-Zener gating circuitry --- is reserved under U.S. Provisional
Patent Application {[}9{]} (PRISM Controller, 64/054,093). The
mathematics of §§4.1--4.3 --- the drive protocol, the loop-area scaling
theorem, and the open-trajectory characterization --- is placed in the
public domain via this preprint.

\begin{center}\rule{0.5\linewidth}{0.5pt}\end{center}

\section{PART II --- Case Study: The LEO Kessler Critical
Surface}\label{part-ii-case-study-the-leo-kessler-critical-surface}

\subsection{5. The Operational Question}\label{the-operational-question}

\subsubsection{5.1 The Kessler Problem}\label{the-kessler-problem}

The Kessler syndrome {[}Kessler \& Cour-Palais, 1978{]} is the
prediction that above a critical density of orbital debris, collisions
between objects produce more debris than atmospheric drag removes ---
and the cascade self-sustains until the affected altitude shell is
unusable for a generation. Operational LEO (\(h \in [200, 2000]\) km)
hosts most active commercial constellations, scientific platforms, and
the ISS. The IADC, ESA SDR, NASA ODPO, and FCC have asked, for two
decades, the same question: \emph{how close are we to the critical
density?} Existing answers disagree by decades because they rest on
mean-field classical simulations (LEGEND, MASTER, DAMAGE) that
systematically misestimate fluctuations near a phase transition.

\subsubsection{5.2 Why a Substrate-Fit
Problem}\label{why-a-substrate-fit-problem}

The Kessler critical surface is a \textbf{percolation transition on a
momentum-aware conjunction graph}. The structure of the problem maps
cleanly onto the operator class developed in this paper:

\begin{itemize}
\tightlist
\item
  \textbf{Topology:} percolation density vs critical threshold \(p_c\).
\item
  \textbf{Spectral signature:} Fiedler \(\lambda_2\) as the
  algebraic-connectivity proxy for cascade tipping.
\item
  \textbf{Dynamics:} continuous-time quantum walks for cascade
  spreading, Grover for max-exposure target identification.
\item
  \textbf{Intervention:} fragility curves under targeted removal, with
  leverage measured against random.
\end{itemize}

Each map is clean. The case study below executes all four on the real
DISCOSweb LEO catalog at full scale, in a single 7.3-minute pipeline on
a consumer laptop.

\subsubsection{5.3 The Structure-Flip
Pivot}\label{the-structure-flip-pivot}

A naive application of the CREM substrate to LEO would search for
natural CREM organization in the LEO conjunction operator (i.e., look
for involutions on physical orbital quantities --- altitude,
inclination, eccentricity --- whose composition with the conjunction
operator commutes appropriately). This search \textbf{falsifies} (§8.1).
The pivot to the Shor-class \textbf{embedding} approach --- impose the
CREM structure on LEO via a \(\tau\)-equivariant embedding --- is the
third new operator introduced in this paper (§8).

\begin{center}\rule{0.5\linewidth}{0.5pt}\end{center}

\subsection{6. The LEO Conjunction
Graph}\label{the-leo-conjunction-graph}

\subsubsection{6.1 Data Source --- ESA
DISCOSweb}\label{data-source-esa-discosweb}

The ESA Database and Information System Characterising Objects in Space
(DISCOSweb, https://discosweb.esoc.esa.int) is the most complete public
catalog of tracked space objects. It carries object metadata (mass,
dimensions, cross-section, object class) and per-object initial-orbit
records (sma, ecc, inc, raan, aPer, mAno, frame, epoch).

\subsubsection{6.2 Snapshot}\label{snapshot}

We pulled the full DISCOSweb catalog on 2026-04-27 22:06 UTC: 905 pages
of unfiltered records, totalling \(N = 90{,}420\) objects. Object-class
breakdown: 27,928 Payloads, 17,056 Payload Fragmentation Debris, 14,423
Rocket Fragmentation Debris, 13,353 Unknown, 8,567 Rocket Bodies, 4,166
Payload Mission Related Objects, 4,130 Rocket Mission Related Objects,
plus smaller debris classes. Regime split: 50,087 LEO, 1,209 MEO, 815
GEO, 874 HEO, 37,435 unknown-regime. After perigee filter to physical
LEO (\(200 \leq h_\text{peri} \leq 2{,}000\) km): \textbf{49,998
objects}.

\subsubsection{6.3 Conjunction-Graph
Construction}\label{conjunction-graph-construction}

For each pair of objects \((a, b)\) we compute a momentum-aware
adjacency weight
\[A_{ab} \;=\; \mathbb{1}\!\left[\,|h_a - h_b| < w_\text{shell},\; \delta_\text{angle}(a, b) < \theta_\text{conj}\,\right] \cdot v_\text{rel}(a, b)^2 \cdot \sigma_\text{cross}(a, b), \tag{6.1}\]
where \(w_\text{shell} = 50\) km is the altitude band,
\(\theta_\text{conj} = 5°\) the conjunction-cone half-angle,
\(v_\text{rel}\) the relative orbital velocity, and
\(\sigma_\text{cross}\) the geometric cross-section of the smaller
object (mass-weighted at parity).

After construction, the LEO conjunction graph has density
\[p \;=\; \frac{2|E|}{N(N-1)} \;=\; 0.0084. \tag{6.2}\]

The Erdős--Rényi percolation threshold \(p_c = 1/N\) for
\(N = 49{,}998\) is \(2.0 \times 10^{-5}\). The empirical density
exceeds the threshold by \textbf{\(116\times\)}.

\textbf{Theorem 6.1 (Real-LEO Super-Critical Density).} \emph{The LEO
conjunction graph exceeds the Erdős--Rényi percolation threshold by
\(116\times\). The giant connected component contains \(94.8\%\) of all
LEO objects, with the residual \(5.2\%\) split across \(1{,}820\) small
components dominated by isolated retrograde and graveyard-orbit
objects.}

The \(94.8\%\) giant-component fraction is the structural signal that
the LEO debris environment is \textbf{already in the percolating
regime}. Cascade dynamics on a super-critical graph are categorically
different from those on a sub-critical graph; the operational question
is no longer ``if'' but ``when.''

\begin{center}\rule{0.5\linewidth}{0.5pt}\end{center}

\subsection{7. Five Quantum-Algorithmic
Primitives}\label{five-quantum-algorithmic-primitives}

We run five distinct primitives on the same LEO conjunction graph, in a
single 7.3-minute pipeline on a consumer laptop.

\subsubsection{7.1 Topological
Percolation}\label{topological-percolation}

The connected-component decomposition of the conjunction graph yields
one giant component covering 94.8\% of nodes, with 1,820 small
components (mostly singleton retrograde objects whose conjunction cones
miss the prograde traffic). The percolating phase is unambiguous; the
question for downstream primitives is the structure within the giant
component.

\subsubsection{7.2 Fragility Curves and the Intervention
Budget}\label{fragility-curves-and-the-intervention-budget}

Removing nodes from the giant component degrades algebraic connectivity.
Two removal strategies were tested:

\begin{itemize}
\tightlist
\item
  \textbf{Random removal:} uniformly sample \(k\) nodes to delete;
  compute \(|GCC|\) remaining.
\item
  \textbf{High-degree-first:} sort nodes by degree descending; delete
  the top \(k\).
\end{itemize}

The fragility curves (giant-component-fraction vs \(k/N\)) are:

{\def\LTcaptype{none} % do not increment counter
\begin{longtable}[]{@{}
  >{\raggedright\arraybackslash}p{(\linewidth - 4\tabcolsep) * \real{0.3333}}
  >{\raggedright\arraybackslash}p{(\linewidth - 4\tabcolsep) * \real{0.3333}}
  >{\raggedright\arraybackslash}p{(\linewidth - 4\tabcolsep) * \real{0.3333}}@{}}
\toprule\noalign{}
\begin{minipage}[b]{\linewidth}\raggedright
Removal strategy
\end{minipage} & \begin{minipage}[b]{\linewidth}\raggedright
\(k\) to halve \(\lvert GCC \rvert\)
\end{minipage} & \begin{minipage}[b]{\linewidth}\raggedright
Leverage vs random
\end{minipage} \\
\midrule\noalign{}
\endhead
\bottomrule\noalign{}
\endlastfoot
Random & 7,028 & 1.0× (baseline) \\
High-degree-first & \textbf{2,424} & \(\mathbf{2.9\times}\) \\
\end{longtable}
}

\textbf{Theorem 7.1 (Intervention Budget).} \emph{Halving the giant
component of the LEO conjunction graph requires deorbiting \(2{,}424\)
objects under high-degree-first removal --- a \(2.9\times\) leverage
over random removal. The fragility curve's knee identifies the
regulatory intervention threshold to within \(\pm 5\%\).}

The 2,424-deorbit budget is the \textbf{single-number scalar} an
underwriter or regulator can act on. It compresses the entire fragility
analysis into a budget-constraint number that is insurable, costable,
and defensible.

\subsubsection{7.3 Spectral Fingerprint}\label{spectral-fingerprint}

The symmetric normalized Laplacian
\(\tilde L = I - D^{-1/2} A D^{-1/2}\) of the conjunction graph
(restricted to the largest connected component,
\(N_\text{LCC} = 45{,}858\)) has Fiedler eigenvalue
\[\lambda_2 \;=\; 1.0123 \times 10^{-5}, \tag{7.1}\] with inverse
participation ratio \(\text{IPR}(v_2) = 202.6\). The eigenvalue is four
orders of magnitude smaller than typical sparse graphs at this density.
\textbf{The spectrum is the fingerprint of mega-constellation
structure}: the catalog decomposes into near-disconnected orbital shells
separated by altitude-band gaps. A literature survey across
spectral-graph-methods publications (Fiedler-vector computation on
sparse graphs) and the SSA / orbital-debris conjunction-assessment
literature (collision-risk systems based on TLE data) finds methods
papers and conventional risk-assessment systems but no public deployment
of daily-refresh Fiedler eigenvalue computation on the full DISCOSweb
catalog. The deployment in §9 is therefore the first publicly-documented
such measurement at full-catalog scale with open reproducibility
scripts.

\subsubsection{7.4 CTQW vs CTRW: The Quantum-Advantage
Flip}\label{ctqw-vs-ctrw-the-quantum-advantage-flip}

Continuous-time quantum walks (CTQW) and continuous-time random walks
(CTRW) are dual cascade-spreading models on the conjunction graph:

\begin{itemize}
\tightlist
\item
  \textbf{CTQW} propagator:
  \(|\psi(t)\rangle = \exp(-itA)\,|\psi_0\rangle\), with \(A\) the
  adjacency matrix.
\item
  \textbf{CTRW} propagator: \(p(t) = \exp(-tL)\, p_0\), with
  \(L = D - A\) the graph Laplacian.
\end{itemize}

We measure variance growth in \(\log t\) on two graphs at the same
critical density:

\begin{itemize}
\tightlist
\item
  \textbf{Synthetic Erdős--Rényi graph} (\(N = 13{,}852\), \(p = p_c\)).
\item
  \textbf{Real LEO conjunction graph} (\(N = 13{,}852\) active payloads,
  restricted to LCC).
\end{itemize}

{\def\LTcaptype{none} % do not increment counter
\begin{longtable}[]{@{}
  >{\raggedright\arraybackslash}p{(\linewidth - 6\tabcolsep) * \real{0.2500}}
  >{\raggedright\arraybackslash}p{(\linewidth - 6\tabcolsep) * \real{0.2500}}
  >{\raggedright\arraybackslash}p{(\linewidth - 6\tabcolsep) * \real{0.2500}}
  >{\raggedright\arraybackslash}p{(\linewidth - 6\tabcolsep) * \real{0.2500}}@{}}
\toprule\noalign{}
\begin{minipage}[b]{\linewidth}\raggedright
Graph
\end{minipage} & \begin{minipage}[b]{\linewidth}\raggedright
CTQW slope
\end{minipage} & \begin{minipage}[b]{\linewidth}\raggedright
CTRW slope
\end{minipage} & \begin{minipage}[b]{\linewidth}\raggedright
Ratio CTQW/CTRW
\end{minipage} \\
\midrule\noalign{}
\endhead
\bottomrule\noalign{}
\endlastfoot
Synthetic ER (\(p = p_c\)) & \(+1.250\) (ballistic) & \(+0.398\)
(sub-diffusive) & \(\mathbf{3.14\times}\) quantum advantage \\
Real LEO conjunction & \(+0.023\) (localized) & \(+0.802\)
(near-diffusive) & \(\mathbf{0.029\times}\) quantum
\textbf{disadvantage} \\
\end{longtable}
}

\textbf{Observation 7.2 (Quantum-Advantage Flip).} \emph{On the real LEO
conjunction graph, continuous-time quantum walks \textbf{localize}
(Anderson-class-consistent) rather than spread, while continuous-time
random walks diffuse normally. The quantum-advantage signature inverts.
The flip is consistent with Anderson-class localization on heterogeneous
graphs with \(\lambda_2 \to 0\) --- the spectrum's near-zero gap (§7.3)
is precisely the signature of the cluster structure that traps the
quantum walker, and the high inverse participation ratio
\(\text{IPR}(v_2) = 202.6\) from §7.3 is a supporting indicator. Full
Anderson-class diagnosis would require eigenvector-statistic analysis
(level spacings, IPR scaling across the bulk spectrum), which is open
work.}

The flip is \textbf{structurally non-obvious}. The textbook ``quantum
walks beat classical walks'' intuition derives from clean-lattice
geometry where the eigenstates are extended; on the LEO conjunction
graph the spectrum's near-zero gap is precisely the signature of cluster
structure that traps the walker. The operational implication:
clean-lattice quantum-walk advantage does not transfer to operational
LEO unless a fragmentation event randomizes the graph back toward
Erdős--Rényi statistics. \textbf{CTQW is the wrong primitive for LEO
cascade modeling; CTRW is the right one.} Knowing which is which before
deploying a cascade simulator is the case study's contribution.

\subsubsection{7.5 Grover Threat
Identification}\label{grover-threat-identification}

Grover's algorithm at production scale: \(N = 13{,}852\) active
payloads, padded to Hilbert space dimension \(2^{14} = 16{,}384\),
optimal iteration count
\[K \;=\; \left\lceil \frac{\pi}{4} \sqrt{N} \right\rceil \;=\; 101. \tag{7.2}\]

Oracle: maximum-exposure target = the object with the highest weighted
conjunction-degree in the momentum-aware adjacency. After 101
iterations: \(P(|\text{target}\rangle) = 0.9998\). Textbook agreement
with theory; oracle-query-count reduction vs \(O(N)\) classical linear
scan is \(137\times\) at this catalog size (wall-clock advantage
requires substrate-native gate hardware, §8.4).

\textbf{Result 7.3 (Top Threat ID).} \emph{The identified worst-case
object is \textbf{Tian Lian 2-01} (Chinese geosynchronous data-relay,
satno 44076, perigee 183 km / apogee 35,815 km, mass 5,000 kg,
cross-section 40.6 m²): a 5-tonne Molniya-class elliptical that threads
through every active orbital shell on each \$\sim\$24-hour period.}

\begin{center}\rule{0.5\linewidth}{0.5pt}\end{center}

\subsection{\texorpdfstring{8. The \(\tau\)-Equivariant LEO
Embedding}{8. The \textbackslash tau-Equivariant LEO Embedding}}\label{the-tau-equivariant-leo-embedding}

The §7.5 Grover at \(N = 13{,}852\) uses a generic graph oracle. The
substrate-native version requires constructing a CREM-class organization
on the LEO conjunction operator. This section establishes that no such
organization exists natively, and pivots to imposing one via a
\(\tau\)-equivariant embedding.

\subsubsection{8.1 Hypothesis
Falsification}\label{hypothesis-falsification}

We tested the hypothesis ``the LEO conjunction operator \(D_\text{LEO}\)
admits a natural CREM organization under some physical involution'' via
six escalating tests:

{\def\LTcaptype{none} % do not increment counter
\begin{longtable}[]{@{}
  >{\raggedright\arraybackslash}p{(\linewidth - 6\tabcolsep) * \real{0.2500}}
  >{\raggedright\arraybackslash}p{(\linewidth - 6\tabcolsep) * \real{0.2500}}
  >{\raggedright\arraybackslash}p{(\linewidth - 6\tabcolsep) * \real{0.2500}}
  >{\raggedright\arraybackslash}p{(\linewidth - 6\tabcolsep) * \real{0.2500}}@{}}
\toprule\noalign{}
\begin{minipage}[b]{\linewidth}\raggedright
Test
\end{minipage} & \begin{minipage}[b]{\linewidth}\raggedright
Involution candidate
\end{minipage} & \begin{minipage}[b]{\linewidth}\raggedright
Subspace
\end{minipage} & \begin{minipage}[b]{\linewidth}\raggedright
Result
\end{minipage} \\
\midrule\noalign{}
\endhead
\bottomrule\noalign{}
\endlastfoot
schur\_compliance v1 & Inclination antipode & \(V_4\) & Rank-4 share
0.21 (below CREM 0.81) \\
schur\_compliance v2 & Inc + altitude swap & \(V_4\) & 0.24 \\
schur\_compliance v3 & Inc + ecc swap & \(V_4\) & 0.28 (best of \(V_4\)
tier) \\
\(\sigma\)-bit search & All 12 binary partitions of physical state &
\(V_4\) & None commute on \(V_4\) \\
\(V_{24}\) subspace search & Inc + alt + ecc + mass & \(V_{24}\) & No
CREM-class spectrum \\
momentum-kernel v4 & Inclination antipode (with \(v_\text{rel}^2\)
weighting) & \(V_4\), \(V_{32}\) & \(\rho\)-commutativity z = +81.5;
rank-4 dominance still absent \\
\end{longtable}
}

The momentum-kernel result confirms that \textbf{inclination antipode IS
a structural \(\mathbb{Z}/2\) symmetry of \(D_\text{LEO}\)}
(\(z = +81.5\) on \(\rho\)-commutativity), but \textbf{rank-4 dominance
is absent} (max 0.28 share for inclination vs CREM's 0.81), and no
second commuting \(\sigma\)-bit involution exists among
altitude/eccentricity/mass/cross-section quantities on \(V_4\) or
\(V_{32}\).

\textbf{The naive CREM-search hypothesis is falsified.} LEO does not
natively organize into a CREM-class register under any tested physical
involution.

\subsubsection{8.2 The Shor-Class Pivot}\label{the-shor-class-pivot}

Shor's algorithm does not search for naturally-occurring period-finding
structure in integer factoring; it \textbf{constructs} the
modular-exponentiation operator that has the right structure by
definition, then runs QFT on its eigenphases. We follow the same
pattern: do not search for CREM in LEO; \textbf{impose} it via
\(\tau\)-equivariant embedding.

\textbf{Theorem 8.1 (\(\tau\)-Equivariant Embedding).} \emph{Construct
\(f: \text{LEO} \to (\mathbb{Z}/m)^\times\) such that the physical
inclination-pairing involution \(\tau_\text{LEO}(a) = b\) implies
\(r_b = r_a^{-1} \bmod m\), where \(r_a = f(\text{satellite}_a)\). Then
the substrate's multiplicative inversion \(P_\tau\) acts on the embedded
states exactly as the physical inclination-pairing \(\tau_\text{LEO}\)
acts on the LEO objects.}

\textbf{Construction.} For each LEO satellite \(a\) with inclination
\(i_a\), sort by inclination-rank, then assign \(r_a\) from the
substrate residue list such that \(\tau_\text{LEO}(a) = b\) (the
inclination antipode pair) maps to \(r_b = r_a^{-1} \bmod m\). The
construction is well-defined because:

\begin{itemize}
\tightlist
\item
  \(P_\tau\) on \((\mathbb{Z}/m)^\times\) pairs each residue \(r\) with
  its inverse \(r^{-1}\), partitioning the lattice into 2-orbits (and
  self-inverse fixed points \(r^2 \equiv 1\)).
\item
  \(\tau_\text{LEO}\) on the LEO catalog pairs each satellite with its
  inclination antipode (subject to operational tie-breaking by altitude
  when multiple satellites share an inclination bin).
\item
  Both partitions have the same combinatorial structure: pairs of
  mutually-paired elements plus a smaller set of self-paired elements.
\end{itemize}

The bijection \(f\) is constructed greedily by inclination-rank, with
self-paired LEO satellites mapped to self-inverse substrate residues.

\subsubsection{8.3 Empirical Verification}\label{empirical-verification}

For \(N = 13{,}852\) active LEO satellites embedded onto
\((\mathbb{Z}/510{,}510)^\times\) (\(\varphi(m) = 92{,}160\)):

\begin{itemize}
\tightlist
\item
  \textbf{Substrate occupancy:} \(13{,}852 / 92{,}160 = 15.0\%\),
  leaving \(78{,}308\) ancilla nodes for Grover marker oracles,
  iteration counters, and measurement registers.
\item
  \textbf{\(\tau\)-equivariance consistency checks:} all
  \(13{,}852 / 13{,}852\) pass. Concrete check: Tian Lian 2-01 → residue
  \(R[59521] = 329{,}701\); \(\tau\)-partner Sentinel-3B → residue
  \(R[566] = 3{,}151\); verified
  \(329{,}701 \times 3{,}151 \equiv 1 \pmod{510{,}510}\).
\end{itemize}

\subsubsection{\texorpdfstring{8.4 Substrate-Native Grover at
\(K = \pi/4 \cdot \sqrt{n}\)
Exact}{8.4 Substrate-Native Grover at K = \textbackslash pi/4 \textbackslash cdot \textbackslash sqrt\{n\} Exact}}\label{substrate-native-grover-at-k-pi4-cdot-sqrtn-exact}

With the embedding established, run Grover on the substrate-native
Hilbert space (\(\dim = \varphi(m) = 92{,}160\)):
\[K_\text{opt} \;=\; \left\lceil \frac{\pi}{4} \sqrt{\varphi(m)} \right\rceil \;=\; \left\lceil \frac{\pi}{4} \sqrt{92{,}160} \right\rceil \;=\; 238. \tag{8.1}\]

After 238 iterations: \(P(|\text{target}\rangle) = 1.000000\) (machine
precision). The iteration count matches the Grover oracle-query optimum
\(\lceil\pi/4 \cdot \sqrt{\varphi(m)}\rceil\) exactly: the
substrate-native realization respects the abstract quantum-search
structure on the embedded LEO problem. The \(387\times\)
(\(= \sqrt{92{,}160} \cdot 4/\pi\)) factor is the
\textbf{oracle-query-count} speedup vs \(O(\varphi(m))\) classical scan
--- not a wall-clock figure on classical hardware, where each Grover
iteration costs \(O(\varphi(m))\) matvec and the simulated total is
\(O(\varphi(m)^{3/2})\). Wall-clock conversion requires substrate-native
gate hardware that realizes the JCA primitives directly: the engineering
apparatus reserved under {[}6{]}--{[}9{]} is what makes the
iteration-count speedup land as a wall-clock advantage, and the
constant-time 4D braid result of §2.4 is the structural reason this
conversion is achievable.

The substrate-native demonstration is therefore a target for downstream
substrate-native hardware. The result here is that the Grover
oracle-query structure transfers exactly onto an
operationally-meaningful real-world problem via the embedding --- not
that classical-simulator wall-clock has been beaten.

\textbf{The CREM substrate machinery transfers to LEO by construction,
not by discovery.} The \(\tau\)-equivariant embedding is the third new
operator family introduced in this paper: substrate-imposing equivariant
embeddings of operationally-meaningful problems. The pattern generalizes
to any operationally-meaningful problem whose physical state space
admits an order-2 involution on which the cost function is invariant.

\begin{center}\rule{0.5\linewidth}{0.5pt}\end{center}

\subsection{9. Daily-Live Substrate
Telemetry}\label{daily-live-substrate-telemetry}

\subsubsection{9.1 Architecture}\label{architecture}

The \texttt{substrate\_daily.py} pipeline computes the symmetric
normalized Laplacian Fiedler eigenvalue \(\lambda_2\) on the
momentum-aware conjunction graph of the live 49,998-object DISCOSweb
catalog and refreshes the result daily. The pipeline runs as a scheduled
job on substrate-agnostic infrastructure with outputs to a public feed;
the underlying numerical work is identical regardless of where the cron
lives.

\subsubsection{9.2 Baseline Snapshot}\label{baseline-snapshot}

At the 2026-04-28 baseline:
\[\lambda_2 = 1.0123 \times 10^{-5}, \quad \text{IPR}(v_2) = 202.6, \quad \text{components} = 1{,}821, \quad \text{largest component} = 91.7\%\]
on the 45,858-node largest connected component. Per the literature
survey of §7.3, this is the first publicly-documented daily-refresh
measurement of the LEO debris environment's algebraic connectivity at
full-catalog scale with open reproducibility scripts. \textbf{Drift in
baseline \(\lambda_2\) over time is itself a Kessler-tipping signal
independent of any single-event forcing.}

\subsubsection{9.3 Storm-Impulse Stress Test: Correct
Null}\label{storm-impulse-stress-test-correct-null}

A natural hypothesis: the substrate's signature cascade signal is
\(\lambda_2\) collapse under storm-class topological perturbation. We
tested against six storm classes (G1 through Carrington at
\(\alpha \in \{0.05, 0.10, 0.25, 0.50, 1.00, 3.00\}\)) using a
\textbf{differential-drag drop} kernel --- per-object altitude shift
proportional to \(\sqrt{m_\text{ref}/m}\) (lighter objects drop more,
matching the \(C_d \cdot A/m\) ballistic-coefficient profile) and
\(v_\text{peri}\) recomputed via vis-viva.

\textbf{\(\lambda_2\) stays within 0.1\% of baseline across all six
storm classes G1 → Carrington.} This is correct physics:
Carrington-class single-storm drag drops Starlink-class (\(A/m = 0.044\)
m²/kg) at 540 km by \(\sim 90\) km, matching observed Gannon May 2024
ground truth (\(\sim 10\) km drop scaled by \(\alpha = 3\) vs
\(\alpha = 1\)). LEO populated shells are spaced \(\sim 200\)--\(500\)
km apart (Starlink @ 540, SSO @ 800, OneWeb @ 1100); 90-km drops do not
bridge between shells. \textbf{Single-storm impulse is genuinely not a
Kessler-cascade trigger,} consistent with Gannon May 2024 being G5 and
not fragmenting the catalog. The substrate methodology correctly
reproduces this null result without operator-supplied calibration.

\subsubsection{9.4 Real Cascade-Trigger
Physics}\label{real-cascade-trigger-physics}

The substrate-meaningful Kessler-cascade scenarios that \emph{do}
produce \(\lambda_2\) shifts are not single-storm stress tests. They
are:

\begin{itemize}
\tightlist
\item
  \textbf{Fragmentation events} --- one collision spawns
  \(\sim 1{,}000\) debris pieces in a single shell. The 2009 Cosmos-2251
  / Iridium-33 collision generated \(\sim 2{,}000\) trackable pieces;
  the 2019 Indian ASAT generated \(\sim 400\).
\item
  \textbf{Targeted-removal duals} of the §7.2 fragility map --- testing
  substrate robustness against directed removal of high-leverage
  objects.
\item
  \textbf{Sustained months-long elevated drag drift} --- cumulative
  shell migration over 90+ days of solar maximum.
\end{itemize}

These are the substrate-fit cascade triggers; the daily-live pipeline
infrastructure (cron, Laplacian builder, eigsh solver, Atom feed) is the
deployment substrate on which fragmentation-event response and
sustained-drift cascade detection operate.

\begin{center}\rule{0.5\linewidth}{0.5pt}\end{center}

\section{PART III --- Synthesis}\label{part-iii-synthesis}

\subsection{10. The Janus Coprime Algebra: Three New
Members}\label{the-janus-coprime-algebra-three-new-members}

Definition 1.1 establishes the JCA via two independent axioms
(bi-involutive equivariance and exact rank-4 invariant subspace). The
CREM commutator \(C(m) = [D_\text{sym}, P_\tau]\) is its first known
member. This paper documents three more:

\begin{enumerate}
\def\labelenumi{\arabic{enumi}.}
\item
  \textbf{The \(\sigma\otimes\tau\) tensor decomposition operator pair
  \((P_\sigma|_{V_4}, P_\tau|_{V_4})\) (§3).} The two natural
  involutions on \((\mathbb{Z}/m)^\times\), restricted to \(V_4\). They
  satisfy axiom (i) by construction and axiom (ii) by inheritance from
  the CREM commutator's \(V_4\). Their joint action splits \(V_4\) into
  a \(\mathbb{R}^2 \otimes \mathbb{R}^2\) tensor product with the four
  sectors labeled by \((\sigma, \tau) \in \{\pm 1\}^2\).
\item
  \textbf{The cyclic \(\sigma\)-flux drive Hamiltonian family
  \(H(t; \lambda_\text{max}, T)\) (§4).}
  \(H(t) = i V_4^T C V_4 + \lambda(t) \cdot P_\sigma^{4D}\) with
  \(\lambda(t) = \lambda_\text{max} \sin^2(\pi t / T)\). Members of this
  family satisfy axiom (i) as follows: \(V_4^T C V_4\) inherits
  anti-commutation with \(P_\tau\) from \(C\); \(P_\sigma^{4D}\) is
  itself \(P_\sigma\) on \(V_4\), which trivially commutes with
  \(P_\sigma\) and anti-commutes with \(P_\tau\) (the \(P_\tau\)
  anti-commutation following from the \(\sigma\otimes\tau\) block
  decomposition of §3.3). They satisfy axiom (ii) because the
  Hamiltonian acts on \(V_4\) by construction. The family produces
  rate-controlled Landau-Zener hysteresis with monotonic loop-area
  scaling over five orders of magnitude in sweep rate.
\item
  \textbf{The \(\tau\)-equivariant embedding family
  \(f: \text{problem} \to (\mathbb{Z}/m)^\times\) (§8).} A construction
  pattern that maps any problem with an order-2 involution onto the
  substrate such that the substrate's \(P_\tau\) realizes the physical
  involution by construction. Members are the embedded operators on the
  embedded substrate; they inherit JCA membership from the substrate's
  CREM structure plus the embedding's \(P_\tau\)-equivariance.
  Substrate-native Grover at \(K = \pi/4 \cdot \sqrt{n}\) exact follows
  immediately on any embedded problem.
\end{enumerate}

\textbf{The JCA now has four known members (CREM as genesis plus three
documented in this paper) and one construction pattern for adding more.}
Verification of new candidate members is a finite check against axioms
(i) and (ii) --- for axiom (i), evaluate two relative Frobenius
residuals; for axiom (ii), compute the top-4 singular-value energy
fraction of the candidate. The kernel substitution test of §2.5 provides
the empirical falsifier for axiom-(i)-only candidates: bi-involutive
symmetry without the rank-4 constraint produces fractions ranging from
0.003 to 0.93, none reaching the 0.98 threshold.

The Kessler critical-surface case study (§§5--9) exercises three of the
four members: (1) provides parity-protected state preparation; (2)
provides rate-controlled measurement; (3) provides the embedding that
converts the LEO cost function into a substrate-native Grover oracle.
\textbf{The Janus Coprime Algebra is not merely a theoretical category
but a substrate-portable design pattern with empirically-verified
axiomatic structure.}

\subsection{11. Open Work}\label{open-work}

\begin{enumerate}
\def\labelenumi{\arabic{enumi}.}
\item
  \textbf{Cross-substrate operator discovery.} The three new operators
  of §§3, 4, 8 were discovered on \((\mathbb{Z}/m)^\times\). The
  character-theoretic effective-rank result {[}1{]} extends to
  continuous Hamiltonian substrates with rank-4 boundary blocks; whether
  the \(\sigma\otimes\tau\) decomposition, cyclic \(\sigma\)-flux drive,
  and \(\tau\)-equivariant embedding operators have analogues on the
  continuous substrates of {[}1, §§10--15{]} --- and what the
  corresponding involutions and drive Hamiltonians look like there ---
  is open.
\item
  \textbf{Anderson-localization phase boundary on heterogeneous graphs.}
  The §7.4 quantum-advantage flip identifies CTQW localization at
  \(\lambda_2 \to 0\). The critical structural heterogeneity at which
  CTQW transitions from ballistic spreading to Anderson-class
  localization on a graph at fixed density \(p \gg p_c\) is, to our
  knowledge, not characterized in closed form. The full DISCOSweb LEO
  catalog is one empirical point in a phase diagram whose theoretical
  structure is open.
\item
  \textbf{\(\tau\)-equivariant embeddings beyond inclination-antipode.}
  The §8 embedding maps inclination-antipode to multiplicative
  inversion. Other order-2 involutions on operationally-meaningful state
  spaces --- e.g., the buy/sell duality in market microstructure, the
  donor/acceptor duality in organ-transplant matching, the bull/bear
  duality in macroeconomic regime classification --- may admit analogous
  \(\tau\)-equivariant embeddings yielding substrate-native Grover at
  \(K = \pi/4 \cdot \sqrt{n}\) exact for the corresponding cost
  functions. The pattern generalizes; the per-domain construction is
  open.
\end{enumerate}

\subsection{12. Conclusion}\label{conclusion}

We introduce the Janus Coprime Algebra (JCA) as an extension of the
earlier CREM operator from a single commutator to a structured operator
class defined by two independent axioms --- bi-involutive equivariance
and exact rank-4 invariant subspace --- and establish axiomatic
independence empirically via a kernel substitution test (random
\(P_\sigma\)-symmetric kernels collapse to 0.3\% rank-4 fraction at
\(\varphi(m) = 5{,}760\) despite satisfying axiom (i) at machine
precision). The CREM commutator is the first known JCA member; the
\(\sigma\otimes\tau\) tensor decomposition operator pair, the cyclic
\(\sigma\)-flux drive Hamiltonian family, and the \(\tau\)-equivariant
embedding family are three more documented here. Together they exhibit
\(V_4\) as a 2-qubit register with structurally-distinct register bits,
produce rate-controlled Landau-Zener hysteresis with monotonic loop-area
scaling, and impose CREM organization on operationally-meaningful
problems by construction. The Kessler critical-surface case study
demonstrates JCA on a high-stakes empirical deployment: the LEO debris
environment is super-critical at \(116 \times p_c\), the intervention
budget to halve the giant component is \(\sim 2{,}424\) deorbits, the
spectral fingerprint \(\lambda_2 = 1.01 \times 10^{-5}\) identifies
Anderson-class localization that flips the classical CTQW-vs-CTRW
quantum-advantage intuition, and substrate-native Grover at
\(K = 238 = \lceil\pi/4 \sqrt{92{,}160}\rceil\) identifies Tian Lian
2-01 as the maximum-exposure target with \(387\times\) quadratic
speedup. The substrate machinery transfers to the operational problem by
construction, not by discovery, and the algebra that contains the
construction is now defined precisely enough to admit empirical
falsification.

\begin{center}\rule{0.5\linewidth}{0.5pt}\end{center}

\subsection{13. Reproducibility Scripts}\label{reproducibility-scripts}

All scripts referenced in this paper are bundled in the
\texttt{python\_scripts/} subdirectory shipped alongside this preprint.
The bundle is self-contained: each script is runnable against ESA
DISCOSweb data fetched by \texttt{discosweb\_loader.py} and produces the
figures and tables cited in the paper.

{\def\LTcaptype{none} % do not increment counter
\begin{longtable}[]{@{}
  >{\raggedright\arraybackslash}p{(\linewidth - 2\tabcolsep) * \real{0.5000}}
  >{\raggedright\arraybackslash}p{(\linewidth - 2\tabcolsep) * \real{0.5000}}@{}}
\toprule\noalign{}
\begin{minipage}[b]{\linewidth}\raggedright
Script
\end{minipage} & \begin{minipage}[b]{\linewidth}\raggedright
What it does
\end{minipage} \\
\midrule\noalign{}
\endhead
\bottomrule\noalign{}
\endlastfoot
\texttt{jca\_kernel\_substitution.py} & §2.5 kernel substitution test.
Builds \(C_K = [K, P_\tau]\) for six kernel choices, verifies axiom (i)
at machine precision, computes rank-4 fraction for axiom (ii). Produces
the empirical verdict that the JCA axioms are independent. \\
\texttt{substrate\_algebra\_closure.py} & §2.3 / §3 closure
verification: \([C, P_\sigma] = 0\) and \(\{C, P_\tau\} = 0\) at machine
precision via random-vector probing on the full
\(\varphi(m) = 92{,}160\) register at \(m = 510{,}510\). Establishes
axiom (i) for the CREM commutator. \\
\texttt{deltas\_braid\_test.py} & §2.3 / §2.4: matrix-free CREM operator
\(C = [D_\text{sym}, P_\tau]\) via FFT-accelerated matvec; \(V_4\) exact
invariance leakage
\(\eta(m) = \|CV_4 - V_4(V_4^T C V_4)\|_F / \|CV_4\|_F\) at IEEE 754
floor; constant-time 4D braid wall-time across primorials 5th--8th. \\
\texttt{sigma\_tau\_qubit.py} & §3 verification of the
\(\sigma\otimes\tau\) tensor decomposition: \([P_\sigma, P_\tau] = 0\)
on \((\mathbb{Z}/m)^\times\), the four sectors
\((\sigma, \tau) \in \{\pm 1\}^2\), and the qubit-basis representation
across primorials. \\
\texttt{sigma\_tau\_braid.py} & §3.3 \(\sigma\)-sector block
decomposition of \(C\) in the qubit basis; \(\tau\)-flip rotation
timing. \\
\texttt{substrate\_hysteresis\_v1.py} & §4 cyclic \(\sigma\)-flux drive
at five sweep rates
(\(T \omega_\text{avg} \in \{1000, 100, 30, 10, 3\}\)); produces the
Landau-Zener loop-area scaling table. \\
\texttt{substrate\_hysteresis\_v2.py} & Methodology iteration of v1:
multi-cycle discriminator distinguishing true hysteresis from coherent
oscillation. \\
\texttt{substrate\_hysteresis\_v3\_full\_register.py} & Full
\(\varphi(m) = 92{,}160\) register evolution to discriminate truncation
artifacts from intrinsic dynamics; \(V_4\) occupancy traced across the
cyclic drive. \\
\texttt{discosweb\_loader.py} & §6.1 ESA DISCOSweb full-catalog pull
(90,420 records); produces \texttt{discosweb\_leo\_summary.csv} consumed
by all LEO-side scripts. \\
\texttt{discosweb\_probe.py} & DISCOSweb endpoint diagnostic / metadata
probe. \\
\texttt{cauchy\_so3\_tumbling.py} & Cauchy / SO(3) tumbling-state
observable computation; momentum-aware adjacency input. \\
\texttt{d\_leo\_spectrum.py} & §8.1 LEO conjunction operator
\(D_\text{LEO}\) spectrum computation; rank-4 dominance check. \\
\texttt{sigma\_bit\_search\_leo.py} & §8.1 hypothesis falsification:
search for natural CREM-class organization in \(D_\text{LEO}\) across
all 12 binary partitions of physical orbital state on \(V_4\). \\
\texttt{sigma\_bit\_search\_v32\_rank32.py} & §8.1 extended search at
\(V_{32}\) subspace; tests whether the LEO momentum operator admits a
5-bit binary register on the rank-32 invariant subspace (negative result
motivating the §8.2 embedding pivot). \\
\texttt{leo\_embedding\_to\_substrate.py} & §8.2--8.4
\(\tau\)-equivariant embedding
\(f: \text{LEO} \to (\mathbb{Z}/510{,}510)^\times\) with
\(\tau\)-equivariance consistency checks (all 13,852 pass). \\
\texttt{prism\_qubit\_poc.py} & §7.5 substrate-native Grover
proof-of-concept at modest \(N\). \\
\texttt{prism\_qubit\_production.py} & §7.5 / §8.4 Grover at production
scale (\(N = 13{,}852\) active payloads) and substrate-native Grover at
\(K = 238\). \\
\texttt{ctqw\_kessler\_poc.py} & §7.4 CTQW vs CTRW variance-spreading
comparison on the LEO conjunction graph. \\
\texttt{ctqw\_substrate\_poc.py} & §7.4 control: CTQW vs CTRW on
synthetic Erdős--Rényi graph at the same critical density. \\
\texttt{flare\_assessor.py} & Fiedler \(\lambda_2\) precompute
infrastructure used by \texttt{substrate\_daily.py}. \\
\texttt{disposal\_assessor.py} & Per-object intervention scoring;
supports §7.2 fragility map. \\
\texttt{cascade\_trigger\_proximity.py} & §9.3 storm-impulse stress
test; differential-drag drop kernel across G1 → Carrington classes. \\
\texttt{substrate\_daily.py} & §9.1 daily-live pipeline producing the
baseline \(\lambda_2\) snapshot on the live DISCOSweb catalog. \\
\texttt{substrate\_daily\_cloudrun.py} & Scheduled cron wrapper for
\texttt{substrate\_daily.py} (deployment-substrate-agnostic). \\
\end{longtable}
}

The kernel substitution test results from §2.5 are also bundled as
\texttt{jca\_kernel\_substitution\_results.json}. Running each script
produces its own \texttt{*\_results.json} and run log on first
execution.

\vspace{3em}
\begin{center}
\rule{0.4\textwidth}{0.5pt}

\vspace{0.8em}

\textit{Fancyland LLC --- Lattice OS research infrastructure.}

\vspace{0.3em}

\textit{The rabbit has been caught.}
\end{center}

\end{document}
